\documentclass[11pt]{article}
\usepackage[a4paper,margin=25mm]{geometry}
\usepackage[T1]{fontenc}
\usepackage{lmodern,amsmath,amssymb,amsthm,booktabs,longtable,array,microtype,xcolor}
\usepackage{xurl}
\usepackage[breaklinks,colorlinks=true,linkcolor=blue!45!black,urlcolor=blue!45!black]{hyperref}
\newtheorem{proposition}{Proposition}[section]
\newtheorem{corollary}[proposition]{Corollary}
\newtheorem{theorem}[proposition]{Theorem}
\newcommand{\FPdim}{\operatorname{FPdim}}
\newcommand{\Rep}{\operatorname{Rep}}
\renewcommand{\Vec}{\operatorname{Vec}}
\newcommand{\Irr}{\operatorname{Irr}}
\newcommand{\C}{\mathbb C}
\newcommand{\Z}{\mathbb Z}
\newcommand{\Q}{\mathbb Q}
\setlength{\emergencystretch}{3em}
\title{Multiplicity-one fusion rings of ranks nine and ten\\\large Evidence, constructions, and exact classification ledger}
\author{Research package prepared for S\'ebastien Palcoux / FusionFrontier\\\small AI-assisted mathematical analysis; independent human review pending}
\date{23 October 2026}
\begin{document}
\maketitle
\begin{abstract}
This package studies the supplied 142 rank-nine and 194 rank-ten
multiplicity-one based fusion rings over the complex numbers. It combines
standard obstructions, a signed-sphericalization reduction, exact constraints
on restrictions of central simple objects, Morita-theoretic arguments, and
explicit categorical constructions. Every accepted outcome is attached to
the original input line, a tensor hash, a proof, and a deterministic replay.
The counts below are generated from the accompanying ledger. An existence
entry means that at least one unitary categorification is constructed; it
does not classify all categorifications of that ring.
\end{abstract}
\begin{center}
\begin{tabular}{lrrrr}
\toprule Rank & Inputs & Excluded over $\C$ & Unitary models & Unresolved \\
\midrule
9 & 142 & 102 & 40 & 0 \\
10 & 194 & 131 & 57 & 6 \\
Total & 336 & 233 & 97 & 6 \\
\bottomrule
\end{tabular}
\end{center}

\newpage
{\small\tableofcontents}
\newpage
\section{Scope and interpretation}
The identifiers $R9\_1\_i$ and $R10\_1\_i$ refer to physical lines $i$ in
the two supplied compressed multiplication-table files. The tensor convention
is $X_iX_j=\sum_kN_{ij}^{k}X_k$, with the unit in position zero.
The input audit checks integrality, nonnegativity, the multiplicity-one bound,
unit, duality, Frobenius reciprocity, all associativity identities, and absence
of duplicate based rings. The supplied tables agree, up to based isomorphism,
with the corresponding AnyonWiki multiplicity-one census snapshot; the
crosswalk is included. This audits the supplied census, rather than repeating
the original enumeration of every possible based ring of these ranks.

An exclusion is a statement about all fusion categories over $\C$. A test
requiring a spherical or pseudo-unitary category is used only when a separate
argument supplies that hypothesis. Weak integrality supplies pseudo-unitarity
by ENO. For the indicated irrational cases, the signed-sphericalization
certificates and an exact character analysis supply a Galois conjugate with
positive categorical dimensions. A unitary model is an existence proof,
not an inference from successful necessary tests.

The machine-readable ledger is authoritative for the current accepted status.
The supplementary directories contain full mathematical proofs and explicit
phase data. Searches that fail to find a model, floating-point optimization
failures, and experimental induction matrices are not treated as exclusions.
The repository and website publish the same accepted ledger, with individual
case pages linking to the full evidence preserved here.

\section{Standard criteria and their scope}
The baseline replay accepts 474 obstruction receipts on 203 distinct input
rings; overlapping receipts are not added to the case count. It verifies each
certificate directly against the original tensor. The following families
are used.
{\small
\begin{longtable}{p{.34\textwidth}p{.57\textwidth}}
\toprule Criterion & Precise role and scope \\\midrule\endhead
Formal-codegree $d$-numbers & General fusion categories. The regular Casimir
has eigenvalues $\dim(E)f_E$; the $d$-number property is unchanged by a nonzero
rational factor. An irreducible polynomial gives an exact divisibility
failure. \\
Cyclotomicity & Applied to commutative based rings, where irreducible
characters are scalar. Exact polynomial and prime witnesses rule out an
abelian splitting field. \\
Drinfeld-center dimension integrality & Applied to integral rings, whose
categorifications would be pseudo-unitary, including a separately checked
noncommutative implementation. \\
Near-group restrictions & Used with their exact
hypotheses, including the divisibility condition for irrational two-value
dimension profiles and the abelian stabilizer condition for
Tambara--Yamagami fusion rules. \\
Excluded fusion subrings & A listed closed based subring carries its own
direct obstruction certificate. Absence from a realization catalogue is
never sufficient. \\
Involution module constraints & Exact nonnegative integer module-matrix
identities attached to the specified invertible object. \\
Two-dimension induction & The required central dimensions cannot be written
in the nonnegative integral dimension semiring with the prescribed unit
multiplicity. \\
Zero-spectrum and one-spectrum criteria & Direct multiplicity-free
pentagon obstructions, valid for general complex fusion categories. \\
\bottomrule
\end{longtable}
}
The report deliberately does not use a unitary Schur-product obstruction,
an unproved universal Isaacs property, or a low-rank modular classification
as an exclusion of arbitrary fusion categories of ranks nine and ten.
The survey and reference scope checks are in \nolinkurl{work/literature/notes.md}.
The 40 subring receipts include two rank-six subrings for which the replay
supplies a direct dimension-lattice witness, so no catalogue inference remains.

\section{A unified central-restriction obstruction}
Suppose $\mathcal C$ is pseudo-unitary. Let $D=\FPdim(\mathcal C)$,
$d_i=\FPdim(X_i)$, and let $F:\mathcal Z(\mathcal C)\to\mathcal C$ be the
forgetful functor with adjoint $I$. For each irreducible representation $E$
of the complexified Grothendieck ring, put $e_E=\dim E$ and let $f_E$ be its
formal codegree. There is a distinct simple central object $A_E$ occurring
in $I(1)$ with multiplicity $e_E$, with
\[
 \beta_E=\dim A_E=D/f_E,\qquad
 v_{E,i}=[F(A_E):X_i],\qquad v_{E,0}=e_E.
\]
Consequently $v_E$ is a nonnegative integer vector, its class is central,
and $\sum_i d_iv_{E,i}=\beta_E$. Let
\[
 B_{ij}=[F I(X_j):X_i]
       =\sum_{a,k}N_{aj}^{k}N_{k,a^*}^{i}.
\]
If the columns of the full restriction matrix are $v_Z$, then
$B=\sum_Zv_Zv_Z^t$. For any chosen subset of known central simples,
\[
 B-\sum_Ev_Ev_E^t
\]
must be entrywise nonnegative and positive semidefinite. These are necessary
conditions, and their use does not assert that a surviving matrix is
categorifiable.

For scalar characters $\chi,\psi$, the symmetry of the central modular
$S$-matrix yields the particularly useful condition
\[
 \beta_\chi\,\chi(v_\psi)=\beta_\psi\,\psi(v_\chi).
\]
For higher-dimensional $E$, the character normalization is
$S_{A_E,Z}=\beta_E\chi_E(FZ)/e_E$; omitting $e_E$ would be incorrect.
The final rank-ten exclusion replay uses only independently verified scalar
characters, so it requires neither a complete character table nor this
higher-dimensional extension. Ostrik's induction theorem
\cite[Theorem~2.13]{Ostrik3} and Deshpande's character formula
\cite[Equations~(35)--(36)]{Deshpande} provide the categorical input.

This single system unifies dimension-lattice, central-support, reciprocity,
and induction-Gram obstructions. A failure at any stage is a rigorous
exclusion after the required pseudo-unitarity reduction. The exact arithmetic
is performed in rational number fields or cyclotomic fields, with no numerical
tolerance in the final decisions.

\subsection{Twist moments in integral dimension sixty}
For a nonunit simple $X$ of dimension $d$, put $J=I(X)$, so $\dim J=Dd$.
The summands $A_E$ of $I(1)$ have twist one. If their forced total dimension
inside $J$ is $M$, the remaining twist eigenvalues lie on the unit circle.
The indicator identities \cite{NgSchauenburg} give
\[
 \operatorname{Tr}(\theta_J)=0,\qquad
 \operatorname{Tr}(\theta_J^2)=D\nu_2(X),\qquad \nu_2(X)\leq1.
\]
Writing real parts of the residual eigenvalues as $x$, their weighted first
moment is $-M$ and their weighted second moment is
$(Dd+D\nu_2(X)-2M)/2$. Cauchy--Schwarz implies
\[
 M\leq\frac{Dd(d+\nu_2(X))}{3d+\nu_2(X)}
 \leq\frac{Dd(d+1)}{3d+1}.
\]
The denominator is positive for $d\geq1$. The last bound is increasing in
$\nu_2(X)$, so it remains valid for non-self-dual simples, whose indicator
is zero. For the four rings $R9\_1\_2,48,84,101$, exact enumeration of
the low central columns leaves 404 candidate frames, and each violates
this bound for a specified simple. The replay checks every branch and every
moment inequality. Two of these four inputs already had baseline exclusions;
only two are new cases in the total.

\section{A sign-associativity criterion and a rank-nine exclusion}

All categories in this section are fusion categories over $\mathbb C$.
Let $R$ be a multiplicity-free fusion ring whose distinguished basis
$b_0=1,b_1,\ldots,b_{r-1}$ is self-dual. An \emph{admissible sign table}
is an associative integral multiplication table $T$ with unit $b_0$,
$T_{ij}^0=\delta_{ij}$, full symmetry of $T_{ij}^k$ in $i,j,k$, and
$|T_{ij}^k|=N_{ij}^k$. A basis-sign change has the form
$T_{ij}^k\mapsto s_i s_j s_k T_{ij}^k$, where $s_0=1$ and
$s_i\in\{1,-1\}$.

\begin{proposition}[Sign-associativity criterion]
If every admissible sign table of $R$ is a basis-sign change of $R$,
then every fusion categorification of $R$ admits a spherical structure.
Consequently, if no formal codegree $D$ of $R$ has $D/f$ an algebraic
integer for every formal codegree $f$, then $R$ has no categorification.
\end{proposition}
\begin{proof}
Let $\mathcal C$ be a hypothetical categorification and take its
canonical sphericalization $\widetilde{\mathcal C}$
\cite[Remark~3.1]{ENO}\cite[Section~5.1]{Ostrik3}.
The two lifts of each simple object differ by the sign object
$\varepsilon$, of spherical dimension $-1$. Every lift of a self-dual
simple is self-dual: its dual has the same nonzero spherical dimension,
whereas the other lift has the opposite dimension.
Choose a unit-preserving set of lifts. In
$K_0(\widetilde{\mathcal C})/(\varepsilon+1)$, their structure
constants form an admissible sign table. Self-duality makes this
quotient commutative, and Frobenius reciprocity gives full symmetry.
Multiplicity one makes each nonzero coefficient equal to $1$ or $-1$.

By hypothesis, tensoring some selected lifts with $\varepsilon$ makes
all nonzero quotient coefficients positive. The product of two such
lifts has exactly the selected lift in every nonzero channel: the
total multiplicity in the two lifts is one and their difference is
$+1$. Hence the selected lifts span a spherical fusion subcategory of
$\widetilde{\mathcal C}$. The forgetful functor restricts to a tensor
equivalence with $\mathcal C$, proving sphericality.

For a spherical categorification, its dimension $D$ is the formal
codegree of its dimension character. Every $D/f$ is an algebraic
integer by \cite[Theorem~2.13]{Ostrik3}. The last assertion follows.
\end{proof}

\begin{corollary}
The supplied ring $R9\_1\_41$, equivalently the AnyonWiki ring
$\mathrm{FR}^{9,0}_{38}$, has no fusion categorification over
$\mathbb C$.
\end{corollary}
\begin{proof}
Use the zero-based basis order in the supplied tensor. There are
$47$ nonzero unordered triples of non-unit labels. Basis-sign changes
allow the following eight coefficients to be set equal to $+1$:
\[
122,\quad 228,\quad 236,\quad 237,\quad245,\quad247,\quad256,\quad334.
\]
Here $ijk$ denotes $T_{ij}^{k}$. Independence of these eight gauge
conditions is an elementary rank-eight calculation over $\mathbb F_2$.
Thirty-six two-term consequences of associativity give the following
table; the remaining nonzero coefficients are all $+1$:
\[
\begin{array}{c|l}
a&335,347,366,446,455,567,777\\
b&336,344,357,456,477,555,667\\
c&337,346,355,444,457,566,677.
\end{array}
\]
The $b_4$ coefficient in $(b_3b_3)b_4=b_3(b_3b_4)$ then reads
\[
a+ab+c=3.
\]
Since $a,b,c\in\{1,-1\}$, this forces $a=b=c=1$.
Thus every admissible sign table is gauge-equivalent to $R$.
The supplementary certificate lists each of the $36$ source
associativity identities and checks every deduction by exact integer
arithmetic. No indicator normalization is used.

The characteristic polynomial of the regular Casimir matrix is
\[
(x-15)^6(x-2)(x^2-120x+1200).
\]
Thus the possible formal codegrees $D$ are $2$, $15$, and
$60\pm20\sqrt6$. The first two fail because $2/15$ and $15/2$
are nonintegral rationals. For the last two,
\[
N_{\mathbb Q(\sqrt6)/\mathbb Q}\!\left(\frac{60\pm20\sqrt6}{15}\right)
=\frac{16}{3},
\]
so the ratios are not algebraic integers. The proposition excludes
general categorification.
\end{proof}

\paragraph{Relation to the literature.}
Vercleyen--Slingerland, \emph{On Low-Rank Multiplicity-Free Complex
Fusion Categories}, Communications in Mathematical Physics
\textbf{407}, 231 (2026), Remark~4.3,
DOI \texttt{10.1007/s00220-026-05705-2}, record the pivotal exclusion
of this ring. The sign-associativity argument above supplies the
missing reduction from general categorification. It is a new
application of sphericalization to this particular tensor, not a
claim that the pivotality conjecture is known in general.

% Bibliography entries expected in the enclosing document:
% ENO: P. Etingof, D. Nikshych, V. Ostrik, On fusion categories,
% Ann. Math. 162 (2005), 581--642, arXiv:math/0203060.
% Ostrik3: V. Ostrik, Pivotal fusion categories of rank 3,
% Moscow Math. J. 15 (2015), 373--396, arXiv:1309.4822.

\section{Oriented signed sphericalization beyond the self-dual case}
For non-self-dual rings, unrestricted permutation symmetry of a sign table
would be wrong. The general checker instead uses the unit, dual reversal,
and the oriented Frobenius-reciprocity symmetries
$T_{ij}^{k}=T_{j,k^*}^{i^*}$ together with dual-compatible basis signs
$s_i=s_{i^*}$. Associativity supplies exact parity equations and then the
remaining integer constraints. The alternative signed tables and their scalar
characters are enumerated completely modulo the recorded gauge choices.

If a sign table is gauge-equivalent to the original positive table, the
chosen lifts close under tensor product in the sphericalization, furnishing
a spherical tensor section. Other sign branches are eliminated only by
their own exact dimension-character and formal-codegree arguments. A
spherical character is promoted to pseudo-unitarity only after an explicit
Galois and pivotal-sign transformation makes every simple dimension positive.
The audit replays 69 oriented receipts, 40 self-dual receipts, 22 alternative
signed-character inventories, and 77 original-character inventories. It
certifies the reductions needed in this report. This is a finite, tensor-specific
argument, not a general solution of the pivotality conjecture.

\section{Eight rank-nine exclusions from central restrictions}
All fusion categories in this section are over $\mathbb C$. The independent
signed-sphericalization certificates show that a categorification of any of
the eight rings considered here has a Galois conjugate with a spherical
structure whose simple dimensions are its Frobenius--Perron dimensions.
It therefore suffices to exclude pseudo-unitary categorifications.

Write $F\colon\mathcal Z(\mathcal C)\to\mathcal C$ for the forgetful
functor and $I$ for its adjoint. Put $D=\operatorname{FPdim}(\mathcal C)$.
For each irreducible representation $E$ of $K_0(\mathcal C)\otimes\mathbb C$,
Ostrik's Theorem~2.13 in \emph{Pivotal fusion categories of rank~3}
(Moscow Math. J. \textbf{15} (2015), 373--396;
\texttt{arXiv:1309.4822}) gives a simple object $A_E$ of the center such that
\[
 \dim A_E=D/f_E,\qquad [F(A_E):\mathbf1]=\dim E.
\]
Here $f_E$ is the formal codegree. The class $[F(A_E)]$ is central in
$K_0(\mathcal C)$, since a half-braiding gives isomorphisms
$F(A_E)\otimes X\cong X\otimes F(A_E)$.
For one-dimensional characters $\chi,\psi$, the modular character formula
and symmetry of the unnormalized $S$-matrix also give
\begin{equation}\label{eq:central-reciprocity}
 \frac D{f_\chi}\,\chi([F(A_\psi)])
 =\frac D{f_\psi}\,\psi([F(A_\chi)]).
\end{equation}
Indeed the character of $K_0(\mathcal Z(\mathcal C))$ obtained by composing
$\chi$ with $F$ is $[B]\mapsto S_{A_\chi,B}/\dim A_\chi$.
This is the correspondence preceding Ostrik's Theorem~2.13.

\subsection{A dimension-lattice obstruction}
Let $L$ be the additive subgroup of $\mathbb R$ generated by the simple
Frobenius--Perron dimensions of $\mathcal C$. Every $D/f_E$ belongs to $L$:
it is the dimension of the object $F(A_E)$. This gives a necessary
condition stronger than algebraic integrality of $D/f_E$.

For each of
\[
 R9\_1\_15,\quad R9\_1\_68,\quad R9\_1\_88,\quad R9\_1\_116,
\]
the dimension multiset is $(1^4,2,a^2,b^2)$, with
$a\simeq3.603875$ and $b\simeq4.493959$; their minimal polynomials are
\[
 a^3-2a^2-8a+8=0,\qquad b^3-4b^2-4b+8=0.
\]
The three numbers $1,a,b$ are rationally independent, and
$D=24+4a+8b$. The formal codegree $8$ occurs. Hence a hypothetical
categorification would satisfy
\[
 3+\tfrac12a+b=\frac D8\in\mathbb Z+\mathbb Za+\mathbb Zb,
\]
contradicting rational independence. The verifier reconstructs the
Frobenius--Perron character, the degree-three field, the independence,
and the codegree directly from each input tensor.

\subsection{A three-character obstruction}
Put $d=2+\sqrt6$, so $d^2=4d+2$. Each of
\[
 R9\_1\_17,\quad R9\_1\_83,\quad R9\_1\_122
\]
has four invertibles, one simple $Y$ of dimension $2$, and four simples
of dimension $d$. Its global dimension is $D=16(1+d)$.
There are two one-dimensional characters $\chi_+,\chi_-$ of codegree
$16$ and one character $\psi$ of codegree $8$ with the following
properties. Both $\chi_\pm$ take value $1$ on every invertible and $-2$
on $Y$. The four dimension-$d$ simples split into two pairs, on which
$\chi_+$ takes values $\sqrt2$ and $-\sqrt2$, respectively;
$\chi_-=-\chi_+$ on these four objects. The character $\psi$ vanishes
on one pair and is nonzero on both members of the other pair.
These assertions are verified as exact character identities against
all fusion products; no inferred character-table completeness is used.

The central objects $A_+,A_-$ corresponding to $\chi_+,\chi_-$ have
dimension $1+d$ and contain the unit with multiplicity one under $F$.
Rational independence of $1,d$ therefore forces
\[
 [F(A_+)]=1+X,\qquad [F(A_-)]=1+Z
\]
for dimension-$d$ simples $X,Z$. Equation~\eqref{eq:central-reciprocity}
for this pair gives $\chi_+(Z)=-\chi_+(X)$, so $X,Z$ lie in opposite
pairs. The central object $A_\psi$ has dimension $2+2d$ and unit
multiplicity one. Thus its restriction has exactly two invertible
summands and two dimension-$d$ summands, counting multiplicities,
and no copy of $Y$. Consequently
\[
 (\chi_++\chi_-)([F(A_\psi)])=4.
\]
Summing \eqref{eq:central-reciprocity} for $(\chi_+,\psi)$ and
$(\chi_-,\psi)$, and using $D/8=2D/16$, gives
\[
 4=2\bigl(2+\psi(X)+\psi(Z)\bigr).
\]
Hence $\psi(X)+\psi(Z)=0$. But exactly one of these two values is
nonzero, since $X,Z$ lie in opposite pairs. This contradiction excludes
all three rings. It uses only three central simple objects and three
characters, with no enumeration of induction matrices.

\subsection{A central-support obstruction}
The remaining ring $R9\_1\_120$ has the same dimension multiset
$(1^4,2,d^4)$ and global dimension $16(1+d)$, but is noncommutative.
It has a one-dimensional character of formal codegree $16$ (its values
on the dimension-$d$ simples are $\pm i\sqrt2$).
The corresponding central simple again restricts to $1+X$ for a single
dimension-$d$ simple $X$. None of the four possible $X$ is central in
the fusion ring. Indeed, in the original zero-based input order,
$X_1,X_3$ fail to commute with $X_2$, while $X_2,X_4$ fail to commute
with $X_1$. Thus $1+X$ cannot be the restriction of a central object.
The certificate records a differing coefficient for each of the four
commutators. This proves the fourth quadratic-family exclusion.

\paragraph{Verification files.}
\nolinkurl{cubic_lattice.py} verifies the four lattice obstructions;
\nolinkurl{quadratic_exact.py} verifies the three-character and
central-support obstructions. Both reconstruct their input tensors
from the supplied archives and bind their receipts to SHA-256 hashes.
Neither final verifier uses floating-point decisions, optimization,
or the exploratory induction-matrix computations.

\section{Ten rank-ten exclusions from central reciprocity}
The rings
\[
R10\_1\_k,\qquad
k\in\{65,89,90,91,120,143,144,147,162,182\},
\]
have no categorification over $\mathbb C$.
The separately verified signed/Galois reduction in \nolinkurl{work/general_signed}
shows that a hypothetical complex categorification of each listed based
ring would give a pseudo-unitary categorification of the same based ring.
The following obstruction rules out that consequence.

Let $\mathcal C$ be pseudo-unitary, $D=\operatorname{FPdim}(\mathcal C)$,
$F:Z(\mathcal C)\to\mathcal C$ the forgetful functor, and $I$ its adjoint.
For each one-dimensional representation $\chi$ of its complexified fusion
ring, let $f_\chi=\sum_i\chi(X_i)\chi(X_i^*)$ be its formal codegree.
There is a distinct simple $A_\chi\subset I(1)$ with
\[
[F(A_\chi):1]=1,\qquad
\dim A_\chi=\beta_\chi=D/f_\chi.
\]
Moreover, for every central simple $Z$,
\[
S_{A_\chi,Z}=\beta_\chi\chi(F(Z)).
\]
These are Ostrik's induction/formal-codegree theorem and the identification
of the corresponding central character (Ostrik, \emph{Pivotal fusion categories of rank 3},
arXiv:1309.4822, Example 2.9 and the discussion preceding
Theorem 2.13; equivalently Deshpande, arXiv:1611.04486v1, equations
(35)--(36)). Symmetry of $S$ gives the necessary equations
\begin{equation}\label{eq:r10reciprocity}
\beta_\chi\chi(v_\psi)=\beta_\psi\psi(v_\chi),
\qquad v_\chi=[F(A_\chi)].
\end{equation}
Only scalar characters are used, also for the eight noncommutative rings;
no noncommutative character recovery or normalization is needed.

Write $N_{ij}^k$ for the supplied based tensor and define its exact
induction Gram matrix by
\[
B_{ij}=\sum_{a,k}N_{aj}^kN_{ka^*}^{i}.
\]
Adjunction gives $B=MM^{\mathsf T}$, where the columns of $M$ are the
restriction vectors of all simples of $Z(\mathcal C)$. Hence any distinct
chosen $v_\chi$ satisfy
\[
 v_\chi\in\mathbb Z_{\ge0}^{10},\quad
(v_\chi)_0=1,\quad d\cdot v_\chi=\beta_\chi,
\quad 0\le(v_\chi)_i\le\min(B_{0i},\lfloor\sqrt{B_{ii}}\rfloor),
\]
and $Q=B-\sum_\chi v_\chi v_\chi^{\mathsf T}$ is entrywise nonnegative
and positive semidefinite. Also $\sum_\chi v_\chi\le B_{0,*}$
coordinatewise, since every selected column has zeroth coordinate one.
The verifier uses only the entrywise and $2\times2$ principal-minor
consequences of positivity. These weak consequences already suffice.

For every listed ring, the dimension vector consists of four $1$'s,
two $2+2\sqrt2$'s and four $3+2\sqrt2$'s, and
$D=96+64\sqrt2$. All scalar characters and codegrees lie in
$\mathbb Q(\sqrt2,i)$. The certificate stores them using the primitive
element $t=\sqrt2+i$, whose minimal polynomial is $t^4-2t^2+9$.
The independent verifier decodes them into the rational basis
$(1,\sqrt2,i,i\sqrt2)$ and checks each of the $100$ multiplicativity
equations for each distinct character. It checks $f_\chi\beta_\chi=D$
exactly. It independently regenerates every possible $v_\chi$ from the
two rational dimension equations and the displayed finite bounds.
It then exhausts equations~\eqref{eq:r10reciprocity}, the Gram bounds,
and first-row inequalities. All ten finite searches are empty, after
only three or five search nodes each. It does not assume that the stored
scalar table is complete, nor does it impose first-row equality.

The machine receipt is \nolinkurl{exclusion_receipts.json}. The initial
searches are \nolinkurl{reciprocal_partial.py} for the eight noncommutative
rings and \nolinkurl{reciprocal.py} for the two commutative rings 89 and 147.
The separate, standard-library-only replay is
\nolinkurl{work/audit_new_models/verify_rank10_reciprocity.py}; its complete
output is \nolinkurl{work/audit_new_models/rank10_reciprocity_audit.json}.
The signed reduction was replayed separately using the verifier in
\nolinkurl{work/general_signed}; it is not part of this standard-library
scalar-character verifier. Thus the
exclusions concern arbitrary complex fusion categories, not just
pseudo-unitary or spherical categories assumed in advance.

\section{Structural exclusions in dimensions forty and one hundred sixty}
\subsection{A Morita obstruction in dimension forty}
The two rings $R10\_1\_72$ and $R10\_1\_127$ have dimension profile
$(1^4,2^4,\sqrt{10}^{\,2})$. Their invertibles form $V_4$. There is an
involution $h$ fixing each dimension-two simple on both sides, with exact
stabilizer $\{1,h\}$, and exchanging the two irrational-dimensional simples.
Some dimension-two simple is not self-dual. These conditions alone give
the following exclusion, without a pivotal assumption.

Let $\mathcal D$ be the integral subcategory, of dimension twenty.
Since $h$ fixes a simple, its pointed associator is trivial in cohomology:
the fixed simple gives a rank-one module category over $\langle h\rangle$.
Thus $A=1+h$ is a connected separable algebra. There are ten simple right
$A$-modules in $\mathcal D$, all of underlying dimension two: two free
orbits of invertibles, and two module structures on each fixed simple.
For each such module $M$, its internal endomorphism object has dimension
two and contains the unit once. It is therefore $1+g$ for an invertible
$g\ne1$ fixing $M$. The underlying stabilizers force $g=h$.
Every module consequently has two simple left $A$-structures, and
$ {}_A\mathcal D_A$ is pointed of dimension twenty.

Write this pointed category as $\Vec(G,\eta)$. The Morita module has
rank ten, so $\mathcal D$ is a group-theoretical category
$\mathcal C(G,\eta,K,\psi)$ with $|K|=2$. The groups of order twenty
are the two abelian groups, $D_{20}$, $\operatorname{Dic}_5$, and
$F_{20}=C_5\rtimes C_4$ with faithful action. A normal order-two $K$
would make every simple invertible, excluding the abelian and dicyclic
possibilities. In $D_{20}$ a nonnormal $K$ is a reflection subgroup;
every relevant double coset is inverse-stable with trivial stabilizer,
so every dimension-two simple is self-dual, independently of $\eta,\psi$.
Hence $G=F_{20}$.

The full Morita dual $ {}_A\mathcal C_A$ has twenty invertibles and two
simples of dimension $\sqrt{10}$. Its noninvertible stabilizer has order
ten and must be abelian by Liptrap's generalized Tambara--Yamagami theorem.
But $F_{20}$ has a unique index-two subgroup, the nonabelian $D_{10}$.
This contradiction proves both exclusions. The proof allows arbitrary
associators throughout; it does not assume that a known untwisted model
is the only possible categorification of the integral subring.

\subsection{Forced commutativity in dimension one hundred sixty}
The two rings $R10\_1\_83$ and $R10\_1\_177$ have dimensions
$(1^2,2^2,5^6)$ and trivial universal grading. Their nontrivial invertible
$g$ fixes the two dimension-two simples and exchanges the six dimension-five
simples in pairs. ENO's categorical Burnside theorem makes a hypothetical
category solvable. In a prime cyclic extension/equivariantization chain,
the final operation cannot be an extension, because the universal grading
is trivial. It must be a $C_2$-equivariantization
$\mathcal C=\mathcal D^{C_2}$.

The sign permutation forces $\mathcal D$ to have dimensions $(1^5,5^3)$.
Its invertibles form $C_5$, act trivially on the three noninvertibles, and
belong to $\mathcal D_{\rm ad}$. The possible adjoint dimensions are
$5+25k$ and must divide eighty; only five and eighty qualify. Dimension
five cannot contain a dimension-five simple in any homogeneous component,
so $\mathcal D_{\rm ad}=\mathcal D$. A second application of the solvable
last-step argument gives a central $\Rep(C_5)$ whose image is exactly the
five invertibles of $\mathcal D$.

An invertible object has at most one half-braiding when the universal grading
is trivial: the ratio of two half-braidings is a tensor automorphism of
the identity, hence a character of that grading. This central $\Rep(C_5)$
is therefore stable under the $C_2$ action. Its equivariantization is a
central Tannakian subcategory of $\mathcal C$ with dimensions $(1,1,2,2)$,
namely $\Rep(D_{10})$. Direct de-equivariantization gives
$\mathcal C=\mathcal E^{D_{10}}$ with $\FPdim\mathcal E=16$.
The orbit-stabilizer dimension formula forces $\mathcal E$ to be pointed:
every nonunit orbit contributes a dimension-five simple and hence has five
invertible underlying objects. The order-five rotation acts freely on all
fifteen nonidentity elements of its group. Characteristic-center and
two-torsion arguments force that group to be $V=C_2^4$.

Identify $V$ with $\mathbb F_{16}$. A rotation is multiplication by an
order-five element $a$, and a reflection can be taken to be $b(x)=x^4$.
The complete presentation of $H^3(V,\C^*)$ has fourteen binary generators.
Their restrictions to cyclic subgroups are the squarefree monomials of
degrees one through three on $V\setminus\{0\}$. They span the even-parity
subspace of its fifteen functions. Rotation invariance leaves four classes,
with constant values on the three length-five orbits and even total parity.
Four explicit normal forms in the certificate cover these classes.

The rotation tensorator can be made strict because
$H^1(C_5,H^2(V,\C^*))=0$. Relative to a reference reflection tensorator,
the commutator of any other reflection tensorator is a rotation-invariant
alternating form. Every such form is
\[
 B_c(x,y)=\operatorname{Tr}_{16/2}(cxy^4),\qquad c\in\mathbb F_4.
\]
It vanishes on $\mathbb F_4\times\mathbb F_4$, the reflection's fixed
space. The reference tensorators are also symmetric there, as checked
exactly. Consequently every allowed tensorator is symmetric on fixed pairs.
This argument covers arbitrary complex actions, not only a bounded search
of roots of unity.

For equivariant objects $X,Y$, the trace of an element $h$ on their tensor
product's fixed homogeneous degree $z$ is
\[
 \sum_{\substack{x+y=z\\hx=x,\ hy=y}}
 J_h(x,y)^{-1}\chi_X(h,x)\chi_Y(h,y).
\]
The established symmetry makes this invariant under exchanging $X,Y$.
Projective stabilizer characters determine equivariant objects, so the
Grothendieck ring must be commutative. Both input tensors are explicitly
noncommutative, proving the exclusions. The full proof and independent
standard-library replay are in \nolinkurl{work/audit_new_models/INT160_AUDIT.md}
and \nolinkurl{verify_int160_audit.py}.

\section{The generalized Tambara--Yamagami family}
Fourteen rank-ten inputs have profile $(1^8,2^2)$. Liptrap's classification
of the corresponding fusion systems supplies necessary and sufficient
associator data. Eleven have explicit unitary data in the package. The
remaining three, $R10\_1\_118$, $R10\_1\_122$, and $R10\_1\_180$, admit no
complex solution of those equations.

For a positive record the phases are powers of a sixteenth root of unity,
with $\tau=1/2$. The verifier checks the defining biderivation and coherence
equations and all eighty-eight recorded nontrivial Fourier blocks for
unitarity. Liptrap's construction then supplies the entire fusion category.
For a negative record, the stabilizer bicharacters are enumerated completely
(sixteen pairs for $V_4$, four for $C_4$). The remaining multiplicative
equations have the form $\prod_jx_j^{A_{ij}}=i^{b_i}$ with arbitrary
$x_j\in\C^*$. For every allowed bicharacter choice the receipt gives an
integer row vector $y$ with
\[
 yA=0,\qquad yb\not\equiv0\pmod4.
\]
Multiplying the equations with these integer exponents would imply
$1=i^{yb}\ne1$. The obstruction therefore applies over all of $\C^*$;
failure of a finite-phase search is not being used as nonexistence evidence.

\section{Complete near-group branch in dimension twenty-four}
\label{sec:ng24-actions}
Consider a fusion category $\mathcal D$ of type $(1^8,4)$ and a categorical
$C_2$ action whose nontrivial element fixes precisely two invertible simples.
Then its equivariantization has based ring $R9\_1\_12$ or $R9\_1\_108$.
More precisely, the former occurs when the invertibles of $\mathcal D$ form
$D_8$, and the latter when they form $Q_8$. In particular this branch cannot
produce $R9\_1\_62$ or $R9\_1\_87$.

This is an exhaustive statement about complex fusion categories and actions.
Indeed dimension $24=2^3\cdot3$ implies solvability by
\cite[Theorem~1.6]{ENOSolvable}, hence weak group-theoreticity and
unitarizability by \cite[Theorem~5.20]{GHR}.
Reutter's Theorems~1--2 and Corollary~3~\cite{Reutter} allow tensor
autoequivalences and their coherence to be transported to unitary ones.
The near-group rules are
$g\rho=\rho=\rho g$ and $\rho^2=\sum_{g\in G}g+2\rho$.
Izumi's classification \cite[Theorems~1.1 and~6.1]{Izumi} shows
$G=D_8$ or $Q_8$, and supplies exactly three
unitary categories for each group. The explicit admissible tuples below
cover all six. Thus a unitary classification is being used only after
unitarizability has been justified.

\subsection{Normal form for all tensor autoequivalences}
We use Izumi, \emph{A Cuntz algebra approach to the classification of
near-group categories}, arXiv:1512.04288v2, Remark~3.6, Definition~4.8,
Theorem~4.9, and Lemma~6.13~\cite{Izumi}. Let $\pi$ be the unique irreducible
two-dimensional representation of $G$, on $K=\mathbb C^2$. Let $j$ be an
antiunitary commuting with $\pi(G)$ and satisfying $j^2=\epsilon$, with
$\epsilon=+1$ for $D_8$ and $\epsilon=-1$ for $Q_8$. For any third root
of unity $\zeta$, the complete categorical data are
\[
(K,V,U_K,j_1,j_2,\chi,l)
=(\mathbb C^2,\pi,\pi,j,\epsilon\zeta j,\chi,0),
\quad
\pi(g)\pi(h)=\chi_h(g)\pi(h)\pi(g).
\]
The generating fusion isometries are
$S_g\in\operatorname{Hom}(g,\rho^2)$ and
$T(v)\in\operatorname{Hom}(\rho,\rho^2)$, $v\in K$.
The left absorption $g\rho=\rho$ is strict in this presentation.
All morphisms are generated by these isometries and their adjoints
(Remark~3.6).

Every tensor autoequivalence is, up to tensor natural isomorphism,
represented by a group automorphism $\phi$ and a unitary matrix $W$ with
\begin{equation}\label{eq:ng-W}
 W\pi(g)W^*=\pi(\phi(g)),\qquad Wj=jW.
\end{equation}
The representative acts by
\begin{equation}\label{eq:ng-gamma}
 \gamma(S_g)=S_{\phi(g)},\qquad
 \gamma(T(v))=T(Wv).
\end{equation}
Here the normal-form assertion includes all tensorators, not just
permutations of simples. To see this, normalize the one-dimensional left
absorption maps $g\rho\to\rho$ under the functor. The tensor-functor
coherence on $(g,h,\rho)$ then normalizes the pointed multiplication
maps as well. More explicitly, if $\lambda_g$ is the tensorator on the
strict left absorption $g\rho=\rho$, the natural-isomorphism component
on $g$ makes $\lambda_g=1$, and the same coherence forces all pointed
tensorators to one. A further scalar choice on $\rho$ fixes $S_1$.
For a functor with tensorator $L$, its normalized trivalent images are
$L_{\rho,\rho}^{-1}F(S_g)$ and $L_{\rho,\rho}^{-1}F(T)$.
The following formulas refer to these normalized images.
Since $S_g$ is obtained by applying the left $g$ action to $S_1$,
its image is now $S_{\phi(g)}$. The action on $K$ is a unitary $W$.
Covariance of the left $G$ action and of the duality defined by $S_1$
gives exactly~\eqref{eq:ng-W}. Remark~3.6 then determines the functor on
all morphisms. This is also the equivalence-of-tuples normal form in
Definition~4.8 and the proof of Theorem~4.9.

Conversely~\eqref{eq:ng-W} preserves $\chi$ and both $j_1,j_2$, and
$l=0$ supplies no further constraint. The reconstruction formulas of
Section 4 therefore give the tensor autoequivalence~\eqref{eq:ng-gamma}.
Explicitly, the formula for $\rho(S_1)$ uses only
$\sum_g S_g$ and $\sum_iT_i j(T_i)$, both covariant under $\gamma$;
$\rho(S_g)=U(g)\rho(S_1)U(g)^*$ is also covariant. Formula (3.12) for
$\rho(T)$ is covariant because $W$ intertwines $\pi$ and commutes with
$j_1,j_2$. Thus this construction supplies genuine functors.

For each $\phi$, Schur's lemma gives $W$ uniquely up to a scalar. The
condition $Wj=jW$ reduces that ambiguity to $\pm W$. Their quotient is
the functor
\[
\delta(S_g)=S_g,\qquad\delta(T)=-T.
\]
It is tensor-isomorphic to the identity: take $\eta_g=1$ and
$\eta_\rho=-1$, extend tensorially, and check naturality on $S_g,T$.
For example $\eta_{\rho^2}\delta(T)=T\eta_\rho=-T$, while the
$S_g$ equation has both sides $S_g$. These generators suffice for all
naturality equations. Consequently there is no additional invisible
autoequivalence or free tensorator parameter at a fixed $\phi$.

Finally $\operatorname{Aut}_\otimes(\operatorname{Id}_{\mathcal D})=1$.
The relation $g\rho=\rho$ forces the scalar on every $g$ to be one,
and the occurrence of $\rho$ in $\rho^2$ forces the scalar on $\rho$
to be one. Thus a square isomorphism for an involutive tensor functor,
when it exists, is unique, and automatically satisfies the $C_2$ action
coherence. In particular there is no independent sign choice that could
alter the duality calculation below.

\subsection{The square of the action determines the large-object duality}
If $\phi^2=1$, then $W^2$ commutes with $\pi(G)$ and is scalar.
The antiunitary condition makes this scalar $+1$ or $-1$.
For $D_8$ choose the real matrices
\[
 \pi(r)=\begin{pmatrix}0&-1\\1&0\end{pmatrix},\quad
 \pi(s)=\begin{pmatrix}1&0\\0&-1\end{pmatrix},\quad j=\text{complex conjugation}.
\]
An automorphism with exactly two fixed group elements is
$r\mapsto r^{-1},s\mapsto r^k s$, $k=1$ or $3$.
Its intertwiner $W$ is a real reflection, so $\operatorname{tr}W=0$
and $W^2=1$. For $k=1$ one may take
$W=2^{-1/2}\left(\begin{smallmatrix}1&1\\1&-1\end{smallmatrix}\right)$.
The action is strict and its square coherence has $\eta_\rho=+1$.

For $Q_8$ use
\[
 \pi(r)=\begin{pmatrix}i&0\\0&-i\end{pmatrix},\quad
 \pi(s)=\begin{pmatrix}0&1\\-1&0\end{pmatrix},\quad
 j=\pi(s)\circ\text{complex conjugation}.
\]
The unitaries commuting with $j$ form $SU(2)$. Every nontrivial
involutive conjugation of $\pi(Q_8)$ lifts to a noncentral $W\in SU(2)$
with $W^2=-1$ and trace zero. In particular this holds for all six
involutions fixing precisely two group elements. Thus $\gamma^2=\delta$,
and the unique square coherence is $\eta_g=1$, $\eta_\rho=-1$.
It is invariant under $\gamma$, proving the action pentagon explicitly.

An equivariant structure on $\rho$ is a scalar $u$ with
$u^2=\eta_\rho$. Since $\gamma(S_1)=S_1$, the lift of $\rho$ is
self-dual exactly when $u^2=1$. Consequently the two lifts are both
self-dual for $D_8$, and form a dual pair for $Q_8$.
This conclusion is independent of the third-root parameter $\zeta$.

\subsection{Independent finite tensor replay}
The pointed part is the strict equivariantization of $\operatorname{Vec}_G$
by $\phi$. Its simple objects are the two characters over each fixed
$g\in\{1,r^2\}$ and one object for each of the three free orbits.
Products are obtained by counting ordered pairs $(g,h)$ with prescribed
product, and projecting the fixed output fibers onto the two $C_2$
characters. This computation allows noncommutative $G$ throughout.
Write $P_+,P_-$ for the two lifts of $\rho$. A fixed invertible
$(g,t)$ multiplies $P_u$ to $P_{tu}$, and each free-orbit simple multiplies
it to $P_++P_-$. For $P_tP_u$, each free invertible orbit occurs once,
and the fixed outputs are
$(1,\epsilon tu)$ and $(r^2,\epsilon tu)$.
Both $P_+$ and $P_-$ occur once because the action on
$\operatorname{Hom}(\rho,\rho^2)$ is $W$ with trace zero.
These rules determine every coefficient.

The standard-library verifier \nolinkurl{verify_ng_actions.py} enumerates
all $8$ automorphisms of $D_8$ and all $24$ of $Q_8$, selecting respectively
the $2$ and $6$ required involutions. For every involution it supplies
an exact matrix $W=A/\sqrt2$, with Gaussian-integer $A$, and checks
unitarity, the intertwining equations, antiunitary compatibility, trace
zero, and $W^2=\epsilon$. It then constructs the complete equivariant
fusion tensor and checks all associativity identities. Every $D_8$
result matches $R9\_1\_12$ and every $Q_8$ result matches
$R9\_1\_108$, coefficient by coefficient, with an explicit
model-to-input permutation. None matches $R9\_1\_62$ or $R9\_1\_87$.
The complete receipt is \nolinkurl{ng_actions.json}.

This argument resolves the entire $(1^8,4)$ de-equivariantization branch;
it does not by itself dispose of other possible de-equivariantization
profiles of a dimension-forty-eight input.

\section{The remaining integral dimension-forty-eight exclusions}
\label{sec:d48-exclusions}
\begin{theorem}
The based rings $R9\_1\_62$ and $R9\_1\_87$ admit no fusion-category
categorification over $\mathbb C$. The same argument distinguishes them from
$R9\_1\_12$ and $R9\_1\_108$, for which unitary constructions are supplied
separately.
\end{theorem}

\paragraph{First reduction.}
All four rings have simple dimensions $(1^4,2^3,4^2)$, invertible group
$V_4$, and trivial universal grading group. The latter assertion is checked
by generating the adjoint subring from all products $XX^*$.
A categorification is solvable by the Etingof--Nikshych--Ostrik
$p^a q^b$ theorem \cite[Theorem~1.6]{ENOSolvable}. In a sequence of cyclic
prime-order extensions and equivariantizations constructing it
(\cite[Proposition~4.4]{ENOSolvable}), the final step cannot be a nontrivial
extension, since its universal grading is trivial. It must therefore be a
$C_p$-equivariantization. Its canonical $\operatorname{Rep}(C_p)$ embeds
centrally, so $p$ divides the order-four invertible group and $p=2$.
Write $\mathcal C=\mathcal D^{C_2}$, and let $e$ be the corresponding sign
object of $\mathcal C$.

The forgetful functor is completely constrained on simple-object orbits.
A pair $\{X,eX\}$ forgets to one simple of dimension $\operatorname{FPdim}(X)$;
a simple fixed by tensoring with $e$ forgets to two distinct simples,
exchanged by the action, of dimension $\operatorname{FPdim}(X)/2$.
The projective stabilizers are cyclic and have trivial complex Schur
multiplier, so no further projective-representation dimension occurs.
The exact three possibilities are as follows; the entries give the
zero-based input label of $e$.
\[
\begin{array}{c|ccc}
\text{input}&(1^8,4)&(1^4,2,4)&(1^4,2^5)\\\hline
R9\_1\_12&7&1&2\\
R9\_1\_62&5&6&7\\
R9\_1\_87&5&4&3\\
R9\_1\_108&5&7&6
\end{array}
\]
The quotient has dimension~24 in each case. The integer verifier
\nolinkurl{verify_d48_exclusions.py} checks these orbits and every dimension
identity directly against the input tensors.

\paragraph{The near-group branch.}
For profile $(1^8,4)$, writing $\rho$ for the unique noninvertible gives
$g\rho=\rho g=\rho$ and
$\rho^2=\sum_{g\in G}g+2\rho$. The acting involution fixes exactly two
invertibles. The complete near-group action theorem of
Section~\ref{sec:ng24-actions} applies: its equivariantization is
$R9\_1\_12$ for $G=D_8$, or $R9\_1\_108$ for $G=Q_8$.
In particular this branch excludes the two rings in the theorem.
This is an exhaustive categorical statement, including all six associator
classes and all tensorators, not merely a test of strict group actions.

\paragraph[The rank-nine quotient branch]{The $(1^4,2^5)$ branch.}
The involution on the quotient fixes exactly two invertibles and one
dimension-two simple $m$. It exchanges the remaining invertibles and pairs
up the remaining four dimension-two simples. The fixed $m$ is self-dual.
The source products of either of its two equivariant lifts force
\[
 m^2=\sum_{g\in\operatorname{Inv}(\mathcal D)}g.
\]
The exact source products used here, in the input labeling, are
$3\cdot3=0+7+8$ for input~12, $1\cdot2=0+5+8$ for input~62,
$6\cdot6=0+5+7$ for input~87, and $1\cdot2=0+5+8$ for input~108.
The last summand in each product forgets to the two nonfixed invertibles.

The quotient ring is necessarily multiplicity-free. For a nonfixed
quotient simple $Z$, its coefficient in the forgetful image of a unique
source simple is one. Thus the tensor-functor identity and the
multiplicity-free source tensor bound every $N_{XY}^Z$ by one.
Invertible output coefficients are bounded by one by rigidity.
The only remaining possible output is $m$. If an input is invertible the
bound is immediate. Otherwise, if $Y\ne m$, Frobenius reciprocity gives
$N_{XY}^{m}=N_{X^*m}^{Y}\le1$, since $Y$ is nonfixed.
For $Y=m$, reciprocity reduces the coefficient to one in $m^2$, which is
already known. This proves the claim without assuming that a quotient of
a multiplicity-free category is generally multiplicity-free.

A direct binary-tensor enumeration now completes the quotient ring,
without consulting a catalog of quotient rings. Label the quotient so that
its dimensions are $(1,1,1,1,2,2,2,2,2)$ and its involution is
$(2\ 3)(5\ 6)(7\ 8)$. The source sign orbits determine the based
restriction matrix $F$, up to this harmless choice of labels.
The verifier \nolinkurl{direct_quotient.py} enumerates every duality involution
compatible with $F$, dimensions, and the action. It identifies tensor
coefficients only under Frobenius reciprocity and the action, then
exhausts all binary assignments satisfying unit, duality, dimension, and
all identities
\[
 \sum_h N_{ij}^{h}(\mathcal C) F_{hk}
 =\sum_{a,b}F_{ia}F_{jb}N_{ab}^{k}(\mathcal D).
\]
For every variable, the positive linear equations also supply an explicit
upper-bound witness at most one, independently confirming that no larger
nonnegative integer coefficient is omitted. Every surviving assignment
is checked against every associativity identity.
There are respectively $8,8,4,4$ possible dualities for sources
$12,62,87,108$; the numbers of linearly feasible completions are
$36,36,4,4$. Exactly four labeled associative completions remain for
each source. Only after this exhaustive completion does the verifier
compare them with $R9\_1\_117$, and all sixteen comparisons give explicit
based isomorphisms. Thus $R9\_1\_117$ is the only possible quotient,
independently of completeness of any input census.

The complete triality-action classification, established in
\nolinkurl{work/resolve_integral/TRIALITY_ACTIONS.md}, now applies.
Every categorification of $R9\_1\_117$ has neutral category
$\operatorname{Rep}(D_8)$ in its canonical $C_3$ grading, and every compatible
$C_2$ action gives only $R9\_1\_12$ or $R9\_1\_108$.
The argument includes the neutral-category restriction, all extension
associator parameters, and the two possible coherence classes.
Consequently this branch also excludes~62 and~87.

\paragraph{The rank-six branch.}
Here write $G$ for the order-four invertible group and $m,n$ for the
unique simples of dimensions~2 and~4. No assumption that $G$ is $V_4$ is
needed at this stage. Dimension, rigidity, and Frobenius reciprocity force
\[
 gm=mg=m,\quad gn=ng=n,\quad
 m^2=\sum_G g,\quad mn=nm=2n,\quad
 n^2=\sum_G g+2m+2n.
\]
In particular this category, denoted $\mathcal B$, has trivial universal
grading. Solvability once again makes it a $C_2$-equivariantization
$\mathcal B=\mathcal E^{C_2}$. The orbit calculation gives
$\operatorname{FPdim}(\mathcal E)=12$, profile $(1^4,2^2)$; the involution
exchanges its two dimension-two simples and transposes two of its
invertibles.

We give the necessary ring reduction explicitly. Let $A$ be the
order-four group of invertibles of $\mathcal E$, and let $X,Y$ be its
dimension-two simples. Duality shows that the left and right actions of
$A$ on $\{X,Y\}$ coincide. If this action were transitive, its stabilizer
$H$ would have order~2 and
$XX^*=\sum_Hh+Z$ for one dimension-two simple $Z$. This expression is
self-dual, so $Z$ is self-dual and hence both $X,Y$ are self-dual.
Choose $g\notin H$; then $Y=gX=Xg$ and $g^2\in H$, whence
$Y^2=g^2X^2=X^2$. An automorphism exchanging $X,Y$ would have to preserve
the unique noninvertible summand $Z$ of this common square, a contradiction.
Thus $A$ fixes both simples, and $XX^*=YY^*=\sum_Aa$.
If both were self-dual, write $XY=aX+bY$; the dimension identity gives
$a+b=2$, whereas $X(XY)=X^2Y=4Y$ forces $a=0$, and
$(XY)Y=XY^2=4X$ then gives an incompatible invertible summand.
Therefore $X^*=Y$, and associativity, or Frobenius reciprocity, gives
\[
 XY=YX=\sum_Aa,\qquad X^2=2Y,\qquad Y^2=2X.
\]
This is the canonical faithful $C_3$ grading with neutral pointed category
on $A$ and one simple in each nonneutral component.

A rank-one homogeneous module forces the neutral associator to be
cohomologically trivial, so the neutral category is $\operatorname{Vec}_A$.
The Brauer--Picard action of a generator of $C_3$ is an isometry
$M\in O(A\oplus A^*,q)$, $q(a,\chi)=\chi(a)$, whose block
$\beta:A^*\to A$ is invertible. These standard extension facts are also
stated and proved in Jordan--Larson,
\cite[Lemma~2.6 and Proposition~3.1]{JordanLarson}. If $A=C_4$, its split quadratic group has exactly four
automorphisms: $I,-I$, coordinate interchange, and its negative. None has
order~3, so this case is impossible. The verifier exhausts all~256 candidate
matrices modulo~4 to check this assertion.

For $A=V_4$, identify $A\oplus A^*$ with $\mathbb F_2^4$, with
$q(a_1,a_2,\chi_1,\chi_2)=a_1\chi_1+a_2\chi_2$ in exponent notation.
The split orthogonal group has~72 elements. Exactly six of them have order~3
and invertible $\beta$. In every one of these six cases,
\begin{equation}\label{eq:d12-no-isotropic-fixed}
 Mv=v,\quad q(v)=0\quad\Longrightarrow\quad v=0.
\end{equation}
This is checked by exhaustive exact enumeration in the same verifier.
Equivalently, the fixed space is zero or is an anisotropic plane.
No choice of scalar associator or of a lift of the action changes this
necessary fixed-vector statement. Jordan--Larson's six-category
classification \cite[Theorem~1.4]{JordanLarson} provides an independent description of the same two matrix
orbits and their three $C_3$ associator twists.

\paragraph{A characteristic central sign.}
Put $\mathcal K=\mathcal E\rtimes C_2$. It is Morita equivalent to
$\mathcal B$ and is faithfully $S_3$-graded with neutral category
$\operatorname{Vec}_{V_4}$, because the involution exchanges the two
nonneutral $C_3$ components. Any central invertible of $\mathcal K$ has
degree in $Z(S_3)=1$; restriction of its half-braiding to the neutral
category gives an $M$-fixed vector in $A\oplus A^*$.
Its twist is precisely $(-1)^{q(v)}$ for that vector $v$. By
\eqref{eq:d12-no-isotropic-fixed}, every central invertible of twist $+1$
has trivial restriction. The kernel of this restriction consists exactly
of grading characters $\operatorname{Hom}(S_3,\mathbb C^*)=C_2$.
Indeed its underlying object is the unit, its half-braiding is trivial
on the neutral component, and its scalar is constant on each indecomposable
homogeneous bimodule component and multiplicative in the grading.
Both kernel objects have twist $+1$. Hence $\mathcal K$, and therefore
$\mathcal B$, has exactly one nontrivial central invertible of twist $+1$.
This proof does not require that any nonzero invariant fermion lift to the
center.

All categories under consideration are weakly integral and therefore have
the canonical pseudo-unitary spherical structure. The braided equivalence
of centers induced by Morita equivalence preserves its twist.
Moreover $U(\mathcal B)=1$, so the forgetful map
$\operatorname{Inv}Z(\mathcal B)\to\operatorname{Inv}(\mathcal B)$ is injective.
The unique nontrivial boson thus has a nontrivial underlying object and
generates a central Tannakian $\operatorname{Rep}(C_2)$.
Every tensor autoequivalence preserves this characteristic central
subcategory.

Return to the original equivariantization
$\mathcal C=\mathcal B^{C_2}$. Equivariantizing this stable central
$\operatorname{Rep}(C_2)$ gives a central Tannakian subcategory of dimension~4
in $\mathcal C$: the equivariantization remains symmetric, and the
composite forgetful functor to $\operatorname{Rep}(C_2)$ and then to
vector spaces is a symmetric fiber functor. It is $\operatorname{Rep}(V_4)$ or
$\operatorname{Rep}(C_4)$. The latter is impossible since
$\operatorname{Inv}(\mathcal C)=V_4$. Thus all four invertible objects of
$\mathcal C$ belong to central $\operatorname{Rep}(V_4)$.
We may therefore choose any of its three signs for de-equivariantization,
in particular the sign in the $(1^8,4)$ column of the displayed table.
The already exhaustive near-group branch forces $R9\_1\_12$ or
$R9\_1\_108$, excluding~62 and~87 in the last remaining branch. This
completes the proof.

\paragraph{Replay and scope.}
The exact finite receipt is
\nolinkurl{work/resolve_rank9_integral/d48_exclusions_receipt.json}.
Run \nolinkurl{python work/resolve_rank9_integral/verify_d48_exclusions.py}
for the dimension identities, all source orbit profiles, the independent
replay of all completed quotient tensors, the $C_4$ obstruction, and all six
quadratic fixed-space checks. The detailed near-group theorem and its
independent matrix/tensor replay are
\nolinkurl{work/resolve_rank9_neargroup/NG_ACTIONS.tex} and
\nolinkurl{verify_ng_actions.py}. The triality action classification is
\nolinkurl{work/resolve_integral/TRIALITY_ACTIONS.md}.
The exhaustive quotient completion and its receipt are
\nolinkurl{work/resolve_rank9_neargroup/direct_quotient.py} and
\nolinkurl{direct_quotient.json}. These categorical reductions and exact
finite enumerations supply the complete exclusion;
no infeasible numerical search or failure to find a model is used.


\section{Unitary constructions}
\subsection{Established models and audited operations}
The baseline contains 84 realized input rings. Seventy-nine have selected
full tensor-model receipts; the Haagerup--Izumi $C_5$ model and four additional
Liptrap-positive inputs complete this count. Standard sources include
pointed and group-representation categories, $SU(2)$ and adjoint $PSU(2)$
categories, Tambara--Yamagami and metaplectic categories, Deligne products,
equivariantizations, group-theoretical categories, and explicit lower-rank
unitary models from Liu--Palcoux--Ren. Every selected receipt includes a
unit-preserving basis permutation and a complete coefficient comparison.

The construction audit checks more than tensor matching for the less standard
operations. It replays 27 associative $C_2$ zestings, eight central-extension
quotients, six cyclic actions on Tambara--Yamagami categories, and the
Ising-swap $C_4$ subextension. In a $C_2$ zesting, a neutral central
involution supplies the two-cocycle, and the fourth-cohomology obstruction
vanishes; the scalar coherence correction is included in the proof. Bases
for which the required central or braided data have not been supplied are
not used for this operation. The eight quotient constructions instead start
from $SU(2)_4$, a specified central extension of $V_4$, and a central boson;
their fusion tensors follow from the actual free-orbit quotient.

\subsection{Explicit coherent pointed actions}
Six new small models are equivariantizations of pointed unitary categories:
\begin{center}
\begin{tabular}{ll}
\toprule Inputs & Construction \\\midrule
$R9\_1\_47$ & $\Vec(C_{12},\omega)^{C_2}$, inversion action \\
$R9\_1\_98$ & $\Vec(C_3\times V_4,\omega)^{C_2}$, inversion and swap \\
$R10\_1\_61,84,109,158$ & $\Vec(\mathbb F_2^3,\omega)^{S_3}$, coordinate permutations \\
\bottomrule
\end{tabular}
\end{center}
For each model the certificate gives all associator exponents $w$, forward
tensorators $J_g$, and compositors $\mu_{g,h}$ modulo eight. The verifier
checks the pointed pentagon and, for every tuple of indices,
\begin{align*}
J_g(a,b)+J_g(ab,c)+w(a,b,c)
 &=w(ga,gb,gc)+J_g(b,c)+J_g(a,bc),\\
\mu_{g,h}(ab)+J_g(ha,hb)+J_h(a,b)
 &=J_{gh}(a,b)+\mu_{g,h}(a)+\mu_{g,h}(b),\\
\mu_{g,h}(ka)+\mu_{gh,k}(a)
 &=\mu_{h,k}(a)+\mu_{g,hk}(a).
\end{align*}
These equalities define coherent categorical actions. Independent code then
constructs all equivariant simple objects as projective matrix actions on
their homogeneous fibers, verifies the action equations, and computes every
fusion coefficient by exact Hom-space characters. Roots of unity of orders
sixteen or forty-eight suffice for these matrix calculations. Thus neither
the existence nor the fusion rules are inferred from a permutation action
on a based ring alone.

The dimension-160 cases $R10\_1\_39$ and $R10\_1\_138$ are likewise
realized as $D_{10}$-equivariantizations of $\Vec(\mathbb F_{16},\omega)$.
One may take $R10\_1\_39=\Rep(\mathbb F_{16}\rtimes D_{10})$, with
action $v\mapsto a^kv^{4^b}$. The second model uses a nontrivial associator
and explicitly stored fourth-root tensorators. The replay checks all 65,536
pointed pentagons, all action diagrams, the projective stabilizer data,
and the 1,000 fusion coefficients of each input.

\subsection{A twisted subgroup-bimodule model}
For $R9\_1\_92$, put $V=\{a\in\mathbb F_2^3:\sum_i a_i=0\}$ and
$G=V\rtimes S_3\cong S_4$. Restrict one of the coherent $S_3$ actions
above to $V$. The crossed-product associator is
\[
 \Omega((a,g),(b,h),(c,k))
 =w(a,gb,ghc)+J_g(b,hc)-\mu_{g,h}(c)\pmod8.
\]
It is verified directly on every quadruple of the 24-element group. Its
restriction to the chosen reflection subgroup $H=C_2$ is trivial. The
category of bimodules over $\bigoplus_{h\in H}h$ in $\Vec(G,\Omega)$
is therefore unitary. The supplied homogeneous bimodules are checked for
left, right, and mixed associativity; the independent reconstruction forms
their balanced tensor products, checks that both actions descend to the
quotient, and identifies all 729 coefficients with $R9\_1\_92$.

\subsection{The triality model of dimension twenty-four}
Nikshych--Riepel, Example 9.3, constructs a nontrivial $C_3$ extension
of $\Rep(D_8)$, integral of dimension 24, with every noninvertible simple
of dimension two. Their definition of nontrivial means that a homogeneous
component is not equivalent to the regular category as a left module.
An invertible simple in that component would give such an equivalence.
Thus both nonneutral components have exactly two dimension-two simples.

Write the neutral simples as $G=V_4$ and $m$, with $m^2=\sum G$.
Every rank-two invertible $\Rep(D_8)$-module has transitive $G$ action.
Indeed the module-category classification gives reflection subgroups $C_2$
or the nontrivial projective class of $D_8$; both are transitive under the
linear characters. The central $C_2$ alternative has a pointed dual and
cannot be an invertible $\Rep(D_8)$-bimodule. The other subgroups have
one, four, or five projective simples, so no case is omitted.

Label the nonneutral simples $X_0,X_1$ and $Y_0,Y_1$ with $X_i^*=Y_i$.
Their left and right actions are nonzero homomorphisms
$\phi,\psi:V_4\to C_2$. Multiplication by $m$ is the all-ones matrix.
Frobenius reciprocity determines the mixed products:
\[
 X_iY_j=\sum_{\phi(g)=i+j}g+m,\qquad
 Y_iX_j=\sum_{\psi(g)=i+j}g+m.
\]
The only unknowns are
$X_iX_j=a_{ij}Y_0+(2-a_{ij})Y_1$, with $a_{ij}\in\{0,1,2\}$;
duality determines $Y_jY_i$. The exact replay checks all
$3\cdot3\cdot3^4=729$ choices, allowing multiplicity two. Precisely six
are associative, all with distinct left and right kernels and $a_{ij}=1$.
Every one is based-isomorphic to $R9\_1\_117$. This identifies the
existence theorem's category with the supplied tensor. The category is
nilpotent and hence unitarizable by Galindo--Hong--Rowell, Theorem 5.20.

\subsection[An actual unitary model for R10-1-21]{An actual unitary model for $R10\_1\_21$}
Let $\mathcal P=PSU(2)_6$ have simples $1,f,x,y=fx$, with
$f^2=1$, $x^2=y^2=1+x+y$, and $xy=f+x+y$.
The invertible $f$ is a transparent fermion and its pointed associator is
trivial. All four simples are self-dual. Write
$\mathcal W=(\mathcal P\boxtimes\mathcal P)^{C_2}$ for the canonical
permutation equivariantization, a unitary braided fusion category.
Its simples are $d_i^\epsilon=(i\boxtimes i,\epsilon)$ and
$p_{ij}=(i\boxtimes j)\oplus(j\boxtimes i)$ for $i<j$.

The object $h=d_f^-$ is an order-two transparent boson: the two fermionic
self-braidings multiply to $+1$, its pointed associator is trivial, and
its equivariance squares to $+1$. Consequently $A=1\oplus h$ is a
connected central commutative separable algebra, and $\mathcal C=\mathcal W_A$
is an actual unitary fusion category. The free-module functor $F$ is tensor.
The $h$-orbits of size two give one simple each; a fixed simple splits
into two simples, since its stabilizer is cyclic of order two. There are
six free orbits and two fixed objects, namely $p_{1f}$ and $p_{xy}$.
Thus $\mathcal C$ has rank ten and dimensions
$1^4,(2+2\sqrt2)^2,(3+2\sqrt2)^4$.

\paragraph{The fixed invertibles, including their tensor squares.}
Put $s=F(d_1^-)$, and let $g,sg$ be the two summands of $F(p_{1f})$.
The pointed subcategory of $\mathcal P\boxtimes\mathcal P$ can be made
strict. For $X=p_{1f}$ write $e_1=1\boxtimes f$, $e_2=f\boxtimes1$.
The equivariant structure exchanges $e_1,e_2$, whereas that of $h$ has
sign $-1$. An equivariant right $h$-action on $X$ is therefore
\[
 r(e_1\otimes h)=i e_2,\qquad r(e_2\otimes h)=-i e_1;
\]
it squares to the identity. Its negative is the other simple module.
The half-braiding $c_{h,X}$ is $-1$ on both components, so the left
$h$-action used in the balanced product is $-r$. Hence $g\otimes_Ag$
is the cokernel of $r\otimes1+1\otimes r$ inside $X\otimes X$.
This operator is Hermitian, so that cokernel identifies with its kernel.
Its two homogeneous lines are
\[
 e_1e_1+e_2e_2,\qquad e_1e_2-e_2e_1.
\]
They have equivariance $+$ and $-$ respectively and underlie $1\oplus h$.
The remaining right action $1\otimes r$ sends the first line to $i$ times
the second and the second to $-i$ times the first; rescaling the second
line by $i$ identifies this with the regular right $A$-module.
Thus $g^2=1$. Multiplication by $s$ exchanges the two module structures,
so the four invertibles form $C_2\times C_2$, with $s^2=g^2=1$.
This calculation is essential: the choice $h=d_f^+$ gives a different
fixed-point resolution and does not give the multiplicity-one tensor here.

\paragraph{All fusion coefficients.}
Choose $a=F(p_{1x})$ and $b=F(d_x^+)$. Then
\[
\begin{array}{c|cccccccccc}
\text{source}&d_1^+&d_1^-&d_f^+&d_f^-&d_x^+&d_x^-&d_y^+&d_y^-&p_{1f}&p_{xy}\\
F&1&s&s&1&b&sb&sb&b&g+sg&gb+sgb
\end{array}
\]
and $F(p_{1y})=F(p_{fx})=ga$, $F(p_{fy})=a$.
Indeed $s$ interchanges the signs on the diagonal simples, and source
fusion gives $(g+sg)b=F(p_{xy})$ and $(g+sg)a=2F(p_{1y})$.
Since $sb\ne b$ and $g^2=1$, these identify a free orbit of four $b$'s
and an orbit of two $a$'s, with $sa=a$.
Ordinary equivariant tensor products give
\[
\begin{split}
b^2&=1+a+ga+b+sb+gb+sgb,\\
a^2&=1+s+a+ga+b+sb,\\
ab&=a+b+sb+gb+sgb.
\end{split}
\]
For example $d_x^+d_x^+=d_1^++d_x^++d_y^++p_{1x}+p_{1y}+p_{xy}$.
The three displayed rules, commutativity, and the invertible action
therefore determine the entire tensor without choosing any unknown
fusion coefficient. They also show all ten simples are self-dual.

The standalone constructor \nolinkurl{model21.py} generates these rules in the displayed model basis. It checks all $10^4$ associativity
identities, and verifies all $10^3$ coefficients against $R10\_1\_21$
under the zero-based model-to-input basis map
\[
(0,9,1,2,3,8,6,7,5,4).
\]
The resulting model receipt is \nolinkurl{model21.json}. Its inverse
input-to-model map is $(0,2,3,4,9,8,6,7,5,1)$. The independent standard-library
verifier \nolinkurl{work/audit_new_models/verify_model21.py} replays the
fixed-point balance, all $196$ source products, the free-module map, and
the complete original-tensor comparison. Its receipt is
\nolinkurl{work/audit_new_models/model21_audit.json}.
The categorical inputs are the standard unitary $PSU(2)_6$ category,
canonical permutation equivariantization, and central Tannakian
condensation (EGNO, Theorems 8.4.9 and 8.23.3). No reconstruction from
a ring-consistent induction matrix is being asserted.

\subsection[Triality equivariantizations]{Triality equivariantizations: $R9\_1\_12$ and $R9\_1\_108$}
\label{sec:triality-equivariant}

\paragraph{Audit status.}
Both constructions below have passed independent categorical and
coefficient audits. The second uses a direct inner modification of the
first action and the vanishing of its action-coherence obstruction.

Put $D=\operatorname{Rep}(D_8)$, where $D_8$ has order eight.  We use its
Tambara--Yamagami presentation
$D=\operatorname{TY}(V_4,\chi,+\tfrac12)$, with alternating bicharacter
$\chi$.  The standard associators are invariant under every permutation of
the three nonidentity elements of $V_4$.  Consequently these permutations,
with identity tensorators, define a coherent unitary $S_3$ action.
The tensor autoequivalence group is $S_3$, while
$\operatorname{BrPic}(D)\cong S_4$; see Nikshych--Riepel,
\emph{Categorical Lagrangian Grassmannians and Brauer--Picard groups of
pointed fusion categories}, \href{https://arxiv.org/abs/1309.5026}{arXiv:1309.5026},
Section~9.1, and Diatlyk--Luo--Wang--Weller,
\emph{Gauging Non-Invertible Symmetries: Topological Interfaces and
Generalized Orbifold Groupoid in 2d QFT},
\href{https://arxiv.org/abs/2311.17044}{arXiv:2311.17044}, Section~3.5,
Equation~(3.30), for the explicit tensor triality.

\paragraph{A coherent Morita transport.}
The three nonregular invertible left $D$-module classes of rank two are the
two reflection-subgroup modules and the projective $D_8$-representation
module for the nontrivial Schur class.  Their simple objects have internal
End objects $1+h+m$, with the three distinct nontrivial choices of $h$.
Tensor triality therefore permutes these three module classes transitively.
Choose the projective module class and realize it by the Azumaya algebra
$A=\operatorname{End}(V)$ in the symmetric category $D$, where $V$ is a real
two-dimensional irreducible projective representation of $D_8$ arising
from its central cover $D_{16}$.  Thus $A=1+h+m$ and
$A\text{-}\operatorname{Bimod}(D)\simeq D$.

Choose the transposition represented by the actual group involution
\[
 a(r)=r^{-1},\qquad a(s)=sr,
 \qquad r^4=s^2=1,\quad srs=r^{-1}.
\]
It fixes the projective module class.  This invariance has a coherent lift:
the nontrivial Schur cocycle of $D_8$ extends to
$D_8\rtimes_a C_2\cong D_{16}$, because the cover $D_{32}\to D_{16}$
restricts to the nontrivial cover $D_{16}\to D_8$ on its normal $D_8$
subgroup.  Equivalently, writing the projective representation as a
rotation $R$ through $\pi/4$ and a reflection $S$, the real reflection
$T=SR^{1/2}$ satisfies
\[
 T^2=1,\qquad TRT=R^{-1},\qquad TST=SR,
\]
and conjugation by $T$ makes $A$ a coherent equivariant algebra.

There is also no ambiguity between this coherent group action and the
strict TY transposition.  A fixed transposition class has two coherent
$C_2$ actions, distinguished by
$H^2(C_2,\operatorname{Aut}_\otimes(\operatorname{Id}_D))
=H^2(C_2,C_2)=C_2$.
For the strict action the two extensions of $m$ have equivariant scalars
$+1,-1$ and are self-dual; in the other class they have scalars $+i,-i$
and form a dual pair.  The group-automorphism action has self-dual
extensions in $\operatorname{Rep}(D_{16})$, so it is coherently equivalent
to the strict restriction.

Form the crossed product $B=D\rtimes S_3$ and its graded Morita dual
$C=A\text{-}\operatorname{Bimod}(B)$.  The neutral component is again $D$.
Its $C_3$ restriction $E$ is nonregular, since triality moves the module
class of $A$; the preceding triality construction identifies its fusion
ring with $R9\_1\_117$.  Coherent $C_2$ equivariance of $A$ supplies an
invertible $t\in C$ with $t^2\cong1$ and trivial restricted associator.
Thus conjugation by $t$ is an actual coherent $C_2$ action on $E$, and
$C=E\rtimes C_2$.  On the neutral category its equivariantization is
$\operatorname{Rep}(D_{16})$.  Indeed, if $q=st$ in
$D_8\rtimes_a C_2$, then $q^2=r$ and $sqs=q^{-1}$.

\paragraph{Duality of the dimension-four objects.}
The subcategories of a faithfully graded fusion category supported on
conjugate subgroups are Morita equivalent: the sum of components on the
corresponding coset is the equivalence bimodule.  Hence all three
order-two subgroup subextensions of $C$ are Morita equivalent to
$D\rtimes C_2$, and therefore to $\operatorname{Rep}(D_{16})$.
The two nonregular ones have dimension profile $(1^4,2^3)$.
Every category in this Morita class is
$\mathcal C(D_{16},1,H,\psi)$.  This profile forces $H=D_{16}$ or a
reflection $C_2$.  Cyclic normal subgroups give pointed duals; a
Klein-four subgroup or an index-two $D_8$ subgroup already gives eight
invertibles, for either available Schur class.  For $H=D_{16}$ the fusion
ring is the representation ring of $D_{16}$.  For a reflection $C_2$,
the three non-normalizer double cosets are inverse-stable with trivial
stabilizer, and the four invertibles form $V_4$.  In both permitted
cases all simples are self-dual.

Write the nonneutral $C_3$ components of $E$ as
$\{X_0,X_1\}$ and $\{Y_0,Y_1\}$, with $X_i^*=Y_i$.
The self-duality just proved for $X_it$ implies $t(X_i)=Y_i$.
Thus both four-dimensional orbit objects in $E^{C_2}$ are self-dual.

\paragraph{Exact tensor identification.}
Label $E$ by
$0,1,2,3,m,X_0,X_1,Y_0,Y_1$, choosing left and right stabilizers
$\{0,1\}$ and $\{0,2\}$ for the $X$ component.
The restriction of the seven neutral simples of
$\operatorname{Rep}(D_{16})$ is
\[
 0,\ 3,\ 0,\ 3,\ m,\ 1+2,\ m,
\]
and the other two restrictions are $X_i+Y_i$.
Let $F$ be this restriction matrix.  Induction is $F^{\mathsf T}$ by
Frobenius reciprocity, and
\[
 \operatorname{Ind}(X)\otimes W
 =\operatorname{Ind}(X\otimes\operatorname{Forget}(W)),\qquad
 W\otimes\operatorname{Ind}(X)
 =\operatorname{Ind}(\operatorname{Forget}(W)\otimes X).
\]
These formulas determine every product outside the neutral subcategory;
inside it we use the ordinary $D_{16}$ representation ring.  The resulting
multiplicity-free tensor agrees with $R9\_1\_12$ under
\[
 (0,3,2,4,8,7,6,1,5).
\]
The independent verifier
\nolinkurl{work/audit_new_models/verify_triality_equivariant.py}
checks all $729$ coefficients, the restriction homomorphism, and the
$D_{16}$ subgroup calculation.  The certificate is
\nolinkurl{work/resolve_small_models/triality_equivariant_certificate.json}.
The construction is unitary throughout; alternatively, unitarizability
follows from Galindo--Hong--Rowell,
\href{https://arxiv.org/abs/1105.5048}{arXiv:1105.5048}, Theorem~5.20.

\paragraph[The second coherent action and R9-1-108]{The second coherent action and $R9\_1\_108$.}
The following direct construction has passed the independent categorical
audit and supplies the second unitary realization.
Let $\alpha$ denote the accepted action and choose $g\in V_4$ with
$g\alpha(g)=h$, the nontrivial fixed invertible.  Set
$\beta=\operatorname{Ad}_g\circ\alpha$.  Then
$\beta^2\simeq\operatorname{Ad}_h$.
For the nonregular triality action, the fixed central invertible in
$\operatorname{Inv}Z(D)$ is $w=(z,h)$, where $z$ is the central involution
of $D_8$.  Its obstruction to a central lift on $h\in E$ lies in
$H^2(C_3,\mathbb C^\times)=0$; hence $\operatorname{Ad}_h\simeq\operatorname{Id}$.
The obstruction to making $\beta$ a coherent $C_2$ action lies in
\[
 H^3(C_2,\operatorname{Aut}_\otimes(\operatorname{Id}_E))
 =H^3(C_2,C_3)=0,
\]
and the corresponding $H^2$ torsor is trivial.  Any correction is a
character of $U(E)=C_3$, so it is the identity on $D$.
Relative to the symmetric central lift of $h$, the lift $w$ differs by
$z$, changing the neutral square coherence by $-1$ on $m$.  The inner
modification changes the orbit pairing to $X_i+Y_{i+1}$.
These obstruction statements use Etingof--Nikshych--Ostrik,
\emph{Fusion categories and homotopy theory},
\href{https://arxiv.org/abs/0909.3140}{arXiv:0909.3140},
Corollary~7.9(i)~\cite{ENOHomotopy},
and Gelaki--Naidu--Nikshych,
\emph{Centers of graded fusion categories},
\href{https://arxiv.org/abs/0905.3117}{arXiv:0905.3117}, Theorem~3.5.
The detailed central-lift and coherence comparison is in
\nolinkurl{work/resolve_integral/TRIALITY_ACTIONS.md}.

The neutral ring is therefore $\operatorname{Rep}(SD_{16})$, since the
corresponding extension has $t^2=z=r^2$, and $q=st$ satisfies
$q^8=s^2=1$ and $sqs=q^3$.
Replacing the neutral group ring accordingly and using this
pairing in the same induction formulas gives exactly $R9\_1\_108$,
with based permutation
\[
 (0,4,6,8,7,1,2,3,5).
\]
Both this coefficient identification and the categorical coherence of the
modification have passed independent audits. No classification of all
possible involutions is required for these two positive witnesses.



\section{Remaining cases and the boundary of this release}

The following inputs retain unresolved categorifiability status. They are not counted as models or exclusions.

\paragraph{Rank 10.} \texttt{R10\_1\_114}, \texttt{R10\_1\_115}, \texttt{R10\_1\_152}, \texttt{R10\_1\_153}, \texttt{R10\_1\_179}, \texttt{R10\_1\_189}.

The six rank-ten cases $R10\_1\_114,115,152,153,179,189$ are explicitly deferred for a later research phase. The five quadratic cases have dimension profile $(1^4,(2+2\sqrt2)^2,(3+2\sqrt2)^4)$, and the last has invertible group $C_5$ and noninvertible dimension $2+\sqrt5$. Their surviving necessary tests do not establish existence. No new classification claim for these six is included in this release.


\section{Reproduction and evidence boundaries}
Start with \nolinkurl{README.md}, then run \nolinkurl{python verify.py} from the
package root. The deterministic replay uses Python 3 with NumPy, SymPy, and
NetworkX; any optional discovery dependencies are separate. Each family
also has a short standalone command in the ledger. Optimized Python
execution is rejected because some inherited verifiers use assertions.

The two tensor digests are explicitly distinguished. One hashes compact JSON
with separators \texttt{(',',':')}; the other hashes the flattened tensor as
unsigned bytes. The original compressed archives have separate file hashes.
The ledger accounts for every one of the 336 physical input lines exactly once
and asserts that no accepted positive and negative outcomes overlap.

The exact calculations check the finite claims in the accompanying proofs.
They do not formally verify the cited category-theoretic theorems in a proof
assistant. Standard categories are grounded in established constructions;
new pointed-action models carry explicit unit-modulus cochains and complete
coherence checks. The triality construction uses an existence theorem and
an exhaustive identification of its possible fusion tensor. These distinct
forms of evidence are identified in the model records.

\appendix
\section{Complete input assignment}
The following compact table is generated from \nolinkurl{ledger.json}.
E means excluded over $\C$, U means a unitary realization exists, and
? means unresolved. The JSON and CSV files give the proof and replay links.
\subsection*{Rank 9}
\begin{longtable}{llll}
\toprule Case / status & Case / status & Case / status & Case / status \\
\midrule\endhead
\texttt{R9\_1\_1}: \textbf{U} & \texttt{R9\_1\_2}: \textbf{E} & \texttt{R9\_1\_3}: \textbf{E} & \texttt{R9\_1\_4}: \textbf{E} \\
\texttt{R9\_1\_5}: \textbf{E} & \texttt{R9\_1\_6}: \textbf{E} & \texttt{R9\_1\_7}: \textbf{E} & \texttt{R9\_1\_8}: \textbf{U} \\
\texttt{R9\_1\_9}: \textbf{U} & \texttt{R9\_1\_10}: \textbf{U} & \texttt{R9\_1\_11}: \textbf{E} & \texttt{R9\_1\_12}: \textbf{U} \\
\texttt{R9\_1\_13}: \textbf{U} & \texttt{R9\_1\_14}: \textbf{E} & \texttt{R9\_1\_15}: \textbf{E} & \texttt{R9\_1\_16}: \textbf{U} \\
\texttt{R9\_1\_17}: \textbf{E} & \texttt{R9\_1\_18}: \textbf{U} & \texttt{R9\_1\_19}: \textbf{E} & \texttt{R9\_1\_20}: \textbf{E} \\
\texttt{R9\_1\_21}: \textbf{E} & \texttt{R9\_1\_22}: \textbf{U} & \texttt{R9\_1\_23}: \textbf{E} & \texttt{R9\_1\_24}: \textbf{E} \\
\texttt{R9\_1\_25}: \textbf{E} & \texttt{R9\_1\_26}: \textbf{U} & \texttt{R9\_1\_27}: \textbf{E} & \texttt{R9\_1\_28}: \textbf{E} \\
\texttt{R9\_1\_29}: \textbf{E} & \texttt{R9\_1\_30}: \textbf{E} & \texttt{R9\_1\_31}: \textbf{E} & \texttt{R9\_1\_32}: \textbf{E} \\
\texttt{R9\_1\_33}: \textbf{E} & \texttt{R9\_1\_34}: \textbf{U} & \texttt{R9\_1\_35}: \textbf{U} & \texttt{R9\_1\_36}: \textbf{U} \\
\texttt{R9\_1\_37}: \textbf{E} & \texttt{R9\_1\_38}: \textbf{E} & \texttt{R9\_1\_39}: \textbf{E} & \texttt{R9\_1\_40}: \textbf{U} \\
\texttt{R9\_1\_41}: \textbf{E} & \texttt{R9\_1\_42}: \textbf{E} & \texttt{R9\_1\_43}: \textbf{E} & \texttt{R9\_1\_44}: \textbf{U} \\
\texttt{R9\_1\_45}: \textbf{U} & \texttt{R9\_1\_46}: \textbf{E} & \texttt{R9\_1\_47}: \textbf{U} & \texttt{R9\_1\_48}: \textbf{E} \\
\texttt{R9\_1\_49}: \textbf{E} & \texttt{R9\_1\_50}: \textbf{E} & \texttt{R9\_1\_51}: \textbf{E} & \texttt{R9\_1\_52}: \textbf{E} \\
\texttt{R9\_1\_53}: \textbf{E} & \texttt{R9\_1\_54}: \textbf{E} & \texttt{R9\_1\_55}: \textbf{E} & \texttt{R9\_1\_56}: \textbf{E} \\
\texttt{R9\_1\_57}: \textbf{U} & \texttt{R9\_1\_58}: \textbf{U} & \texttt{R9\_1\_59}: \textbf{E} & \texttt{R9\_1\_60}: \textbf{U} \\
\texttt{R9\_1\_61}: \textbf{E} & \texttt{R9\_1\_62}: \textbf{E} & \texttt{R9\_1\_63}: \textbf{U} & \texttt{R9\_1\_64}: \textbf{E} \\
\texttt{R9\_1\_65}: \textbf{E} & \texttt{R9\_1\_66}: \textbf{E} & \texttt{R9\_1\_67}: \textbf{U} & \texttt{R9\_1\_68}: \textbf{E} \\
\texttt{R9\_1\_69}: \textbf{E} & \texttt{R9\_1\_70}: \textbf{U} & \texttt{R9\_1\_71}: \textbf{E} & \texttt{R9\_1\_72}: \textbf{U} \\
\texttt{R9\_1\_73}: \textbf{E} & \texttt{R9\_1\_74}: \textbf{E} & \texttt{R9\_1\_75}: \textbf{E} & \texttt{R9\_1\_76}: \textbf{E} \\
\texttt{R9\_1\_77}: \textbf{E} & \texttt{R9\_1\_78}: \textbf{E} & \texttt{R9\_1\_79}: \textbf{E} & \texttt{R9\_1\_80}: \textbf{E} \\
\texttt{R9\_1\_81}: \textbf{E} & \texttt{R9\_1\_82}: \textbf{E} & \texttt{R9\_1\_83}: \textbf{E} & \texttt{R9\_1\_84}: \textbf{E} \\
\texttt{R9\_1\_85}: \textbf{E} & \texttt{R9\_1\_86}: \textbf{E} & \texttt{R9\_1\_87}: \textbf{E} & \texttt{R9\_1\_88}: \textbf{E} \\
\texttt{R9\_1\_89}: \textbf{E} & \texttt{R9\_1\_90}: \textbf{E} & \texttt{R9\_1\_91}: \textbf{E} & \texttt{R9\_1\_92}: \textbf{U} \\
\texttt{R9\_1\_93}: \textbf{E} & \texttt{R9\_1\_94}: \textbf{U} & \texttt{R9\_1\_95}: \textbf{E} & \texttt{R9\_1\_96}: \textbf{E} \\
\texttt{R9\_1\_97}: \textbf{E} & \texttt{R9\_1\_98}: \textbf{U} & \texttt{R9\_1\_99}: \textbf{E} & \texttt{R9\_1\_100}: \textbf{E} \\
\texttt{R9\_1\_101}: \textbf{E} & \texttt{R9\_1\_102}: \textbf{U} & \texttt{R9\_1\_103}: \textbf{U} & \texttt{R9\_1\_104}: \textbf{U} \\
\texttt{R9\_1\_105}: \textbf{E} & \texttt{R9\_1\_106}: \textbf{E} & \texttt{R9\_1\_107}: \textbf{U} & \texttt{R9\_1\_108}: \textbf{U} \\
\texttt{R9\_1\_109}: \textbf{E} & \texttt{R9\_1\_110}: \textbf{E} & \texttt{R9\_1\_111}: \textbf{U} & \texttt{R9\_1\_112}: \textbf{E} \\
\texttt{R9\_1\_113}: \textbf{E} & \texttt{R9\_1\_114}: \textbf{E} & \texttt{R9\_1\_115}: \textbf{E} & \texttt{R9\_1\_116}: \textbf{E} \\
\texttt{R9\_1\_117}: \textbf{U} & \texttt{R9\_1\_118}: \textbf{E} & \texttt{R9\_1\_119}: \textbf{E} & \texttt{R9\_1\_120}: \textbf{E} \\
\texttt{R9\_1\_121}: \textbf{E} & \texttt{R9\_1\_122}: \textbf{E} & \texttt{R9\_1\_123}: \textbf{E} & \texttt{R9\_1\_124}: \textbf{E} \\
\texttt{R9\_1\_125}: \textbf{E} & \texttt{R9\_1\_126}: \textbf{E} & \texttt{R9\_1\_127}: \textbf{E} & \texttt{R9\_1\_128}: \textbf{U} \\
\texttt{R9\_1\_129}: \textbf{E} & \texttt{R9\_1\_130}: \textbf{E} & \texttt{R9\_1\_131}: \textbf{U} & \texttt{R9\_1\_132}: \textbf{U} \\
\texttt{R9\_1\_133}: \textbf{E} & \texttt{R9\_1\_134}: \textbf{E} & \texttt{R9\_1\_135}: \textbf{E} & \texttt{R9\_1\_136}: \textbf{U} \\
\texttt{R9\_1\_137}: \textbf{E} & \texttt{R9\_1\_138}: \textbf{E} & \texttt{R9\_1\_139}: \textbf{E} & \texttt{R9\_1\_140}: \textbf{E} \\
\texttt{R9\_1\_141}: \textbf{U} & \texttt{R9\_1\_142}: \textbf{U} &  &  \\
\bottomrule
\end{longtable}
\subsection*{Rank 10}
\begin{longtable}{llll}
\toprule Case / status & Case / status & Case / status & Case / status \\
\midrule\endhead
\texttt{R10\_1\_1}: \textbf{U} & \texttt{R10\_1\_2}: \textbf{E} & \texttt{R10\_1\_3}: \textbf{E} & \texttt{R10\_1\_4}: \textbf{U} \\
\texttt{R10\_1\_5}: \textbf{E} & \texttt{R10\_1\_6}: \textbf{E} & \texttt{R10\_1\_7}: \textbf{E} & \texttt{R10\_1\_8}: \textbf{U} \\
\texttt{R10\_1\_9}: \textbf{E} & \texttt{R10\_1\_10}: \textbf{U} & \texttt{R10\_1\_11}: \textbf{E} & \texttt{R10\_1\_12}: \textbf{E} \\
\texttt{R10\_1\_13}: \textbf{E} & \texttt{R10\_1\_14}: \textbf{U} & \texttt{R10\_1\_15}: \textbf{E} & \texttt{R10\_1\_16}: \textbf{E} \\
\texttt{R10\_1\_17}: \textbf{U} & \texttt{R10\_1\_18}: \textbf{E} & \texttt{R10\_1\_19}: \textbf{E} & \texttt{R10\_1\_20}: \textbf{E} \\
\texttt{R10\_1\_21}: \textbf{U} & \texttt{R10\_1\_22}: \textbf{U} & \texttt{R10\_1\_23}: \textbf{E} & \texttt{R10\_1\_24}: \textbf{U} \\
\texttt{R10\_1\_25}: \textbf{E} & \texttt{R10\_1\_26}: \textbf{E} & \texttt{R10\_1\_27}: \textbf{E} & \texttt{R10\_1\_28}: \textbf{E} \\
\texttt{R10\_1\_29}: \textbf{E} & \texttt{R10\_1\_30}: \textbf{E} & \texttt{R10\_1\_31}: \textbf{U} & \texttt{R10\_1\_32}: \textbf{E} \\
\texttt{R10\_1\_33}: \textbf{U} & \texttt{R10\_1\_34}: \textbf{E} & \texttt{R10\_1\_35}: \textbf{U} & \texttt{R10\_1\_36}: \textbf{U} \\
\texttt{R10\_1\_37}: \textbf{E} & \texttt{R10\_1\_38}: \textbf{E} & \texttt{R10\_1\_39}: \textbf{U} & \texttt{R10\_1\_40}: \textbf{U} \\
\texttt{R10\_1\_41}: \textbf{U} & \texttt{R10\_1\_42}: \textbf{E} & \texttt{R10\_1\_43}: \textbf{E} & \texttt{R10\_1\_44}: \textbf{E} \\
\texttt{R10\_1\_45}: \textbf{E} & \texttt{R10\_1\_46}: \textbf{E} & \texttt{R10\_1\_47}: \textbf{U} & \texttt{R10\_1\_48}: \textbf{E} \\
\texttt{R10\_1\_49}: \textbf{E} & \texttt{R10\_1\_50}: \textbf{E} & \texttt{R10\_1\_51}: \textbf{U} & \texttt{R10\_1\_52}: \textbf{U} \\
\texttt{R10\_1\_53}: \textbf{E} & \texttt{R10\_1\_54}: \textbf{U} & \texttt{R10\_1\_55}: \textbf{E} & \texttt{R10\_1\_56}: \textbf{E} \\
\texttt{R10\_1\_57}: \textbf{U} & \texttt{R10\_1\_58}: \textbf{E} & \texttt{R10\_1\_59}: \textbf{E} & \texttt{R10\_1\_60}: \textbf{U} \\
\texttt{R10\_1\_61}: \textbf{U} & \texttt{R10\_1\_62}: \textbf{E} & \texttt{R10\_1\_63}: \textbf{E} & \texttt{R10\_1\_64}: \textbf{E} \\
\texttt{R10\_1\_65}: \textbf{E} & \texttt{R10\_1\_66}: \textbf{E} & \texttt{R10\_1\_67}: \textbf{E} & \texttt{R10\_1\_68}: \textbf{U} \\
\texttt{R10\_1\_69}: \textbf{E} & \texttt{R10\_1\_70}: \textbf{E} & \texttt{R10\_1\_71}: \textbf{U} & \texttt{R10\_1\_72}: \textbf{E} \\
\texttt{R10\_1\_73}: \textbf{E} & \texttt{R10\_1\_74}: \textbf{E} & \texttt{R10\_1\_75}: \textbf{E} & \texttt{R10\_1\_76}: \textbf{U} \\
\texttt{R10\_1\_77}: \textbf{E} & \texttt{R10\_1\_78}: \textbf{E} & \texttt{R10\_1\_79}: \textbf{U} & \texttt{R10\_1\_80}: \textbf{U} \\
\texttt{R10\_1\_81}: \textbf{E} & \texttt{R10\_1\_82}: \textbf{U} & \texttt{R10\_1\_83}: \textbf{E} & \texttt{R10\_1\_84}: \textbf{U} \\
\texttt{R10\_1\_85}: \textbf{E} & \texttt{R10\_1\_86}: \textbf{E} & \texttt{R10\_1\_87}: \textbf{E} & \texttt{R10\_1\_88}: \textbf{E} \\
\texttt{R10\_1\_89}: \textbf{E} & \texttt{R10\_1\_90}: \textbf{E} & \texttt{R10\_1\_91}: \textbf{E} & \texttt{R10\_1\_92}: \textbf{E} \\
\texttt{R10\_1\_93}: \textbf{E} & \texttt{R10\_1\_94}: \textbf{E} & \texttt{R10\_1\_95}: \textbf{E} & \texttt{R10\_1\_96}: \textbf{E} \\
\texttt{R10\_1\_97}: \textbf{U} & \texttt{R10\_1\_98}: \textbf{E} & \texttt{R10\_1\_99}: \textbf{E} & \texttt{R10\_1\_100}: \textbf{U} \\
\texttt{R10\_1\_101}: \textbf{E} & \texttt{R10\_1\_102}: \textbf{E} & \texttt{R10\_1\_103}: \textbf{U} & \texttt{R10\_1\_104}: \textbf{E} \\
\texttt{R10\_1\_105}: \textbf{U} & \texttt{R10\_1\_106}: \textbf{U} & \texttt{R10\_1\_107}: \textbf{E} & \texttt{R10\_1\_108}: \textbf{E} \\
\texttt{R10\_1\_109}: \textbf{U} & \texttt{R10\_1\_110}: \textbf{E} & \texttt{R10\_1\_111}: \textbf{E} & \texttt{R10\_1\_112}: \textbf{E} \\
\texttt{R10\_1\_113}: \textbf{E} & \texttt{R10\_1\_114}: \textbf{?} & \texttt{R10\_1\_115}: \textbf{?} & \texttt{R10\_1\_116}: \textbf{U} \\
\texttt{R10\_1\_117}: \textbf{E} & \texttt{R10\_1\_118}: \textbf{E} & \texttt{R10\_1\_119}: \textbf{E} & \texttt{R10\_1\_120}: \textbf{E} \\
\texttt{R10\_1\_121}: \textbf{E} & \texttt{R10\_1\_122}: \textbf{E} & \texttt{R10\_1\_123}: \textbf{E} & \texttt{R10\_1\_124}: \textbf{U} \\
\texttt{R10\_1\_125}: \textbf{U} & \texttt{R10\_1\_126}: \textbf{U} & \texttt{R10\_1\_127}: \textbf{E} & \texttt{R10\_1\_128}: \textbf{E} \\
\texttt{R10\_1\_129}: \textbf{E} & \texttt{R10\_1\_130}: \textbf{E} & \texttt{R10\_1\_131}: \textbf{U} & \texttt{R10\_1\_132}: \textbf{U} \\
\texttt{R10\_1\_133}: \textbf{U} & \texttt{R10\_1\_134}: \textbf{E} & \texttt{R10\_1\_135}: \textbf{E} & \texttt{R10\_1\_136}: \textbf{E} \\
\texttt{R10\_1\_137}: \textbf{E} & \texttt{R10\_1\_138}: \textbf{U} & \texttt{R10\_1\_139}: \textbf{E} & \texttt{R10\_1\_140}: \textbf{U} \\
\texttt{R10\_1\_141}: \textbf{E} & \texttt{R10\_1\_142}: \textbf{E} & \texttt{R10\_1\_143}: \textbf{E} & \texttt{R10\_1\_144}: \textbf{E} \\
\texttt{R10\_1\_145}: \textbf{E} & \texttt{R10\_1\_146}: \textbf{E} & \texttt{R10\_1\_147}: \textbf{E} & \texttt{R10\_1\_148}: \textbf{U} \\
\texttt{R10\_1\_149}: \textbf{E} & \texttt{R10\_1\_150}: \textbf{U} & \texttt{R10\_1\_151}: \textbf{E} & \texttt{R10\_1\_152}: \textbf{?} \\
\texttt{R10\_1\_153}: \textbf{?} & \texttt{R10\_1\_154}: \textbf{U} & \texttt{R10\_1\_155}: \textbf{E} & \texttt{R10\_1\_156}: \textbf{E} \\
\texttt{R10\_1\_157}: \textbf{U} & \texttt{R10\_1\_158}: \textbf{U} & \texttt{R10\_1\_159}: \textbf{E} & \texttt{R10\_1\_160}: \textbf{E} \\
\texttt{R10\_1\_161}: \textbf{E} & \texttt{R10\_1\_162}: \textbf{E} & \texttt{R10\_1\_163}: \textbf{U} & \texttt{R10\_1\_164}: \textbf{E} \\
\texttt{R10\_1\_165}: \textbf{E} & \texttt{R10\_1\_166}: \textbf{E} & \texttt{R10\_1\_167}: \textbf{U} & \texttt{R10\_1\_168}: \textbf{E} \\
\texttt{R10\_1\_169}: \textbf{E} & \texttt{R10\_1\_170}: \textbf{E} & \texttt{R10\_1\_171}: \textbf{E} & \texttt{R10\_1\_172}: \textbf{E} \\
\texttt{R10\_1\_173}: \textbf{E} & \texttt{R10\_1\_174}: \textbf{E} & \texttt{R10\_1\_175}: \textbf{E} & \texttt{R10\_1\_176}: \textbf{E} \\
\texttt{R10\_1\_177}: \textbf{E} & \texttt{R10\_1\_178}: \textbf{E} & \texttt{R10\_1\_179}: \textbf{?} & \texttt{R10\_1\_180}: \textbf{E} \\
\texttt{R10\_1\_181}: \textbf{E} & \texttt{R10\_1\_182}: \textbf{E} & \texttt{R10\_1\_183}: \textbf{U} & \texttt{R10\_1\_184}: \textbf{E} \\
\texttt{R10\_1\_185}: \textbf{U} & \texttt{R10\_1\_186}: \textbf{U} & \texttt{R10\_1\_187}: \textbf{E} & \texttt{R10\_1\_188}: \textbf{E} \\
\texttt{R10\_1\_189}: \textbf{?} & \texttt{R10\_1\_190}: \textbf{E} & \texttt{R10\_1\_191}: \textbf{U} & \texttt{R10\_1\_192}: \textbf{E} \\
\texttt{R10\_1\_193}: \textbf{U} & \texttt{R10\_1\_194}: \textbf{E} &  &  \\
\bottomrule
\end{longtable}


\begin{thebibliography}{99}
\bibitem{ENO} P. Etingof, D. Nikshych, V. Ostrik,
\emph{On fusion categories}, Ann. Math. 162 (2005), 581--642,
\href{https://arxiv.org/abs/math/0203060}{arXiv:math/0203060}.
\bibitem{OstrikCodegrees} V. Ostrik, \emph{On formal codegrees of fusion
categories}, Math. Res. Lett. 16 (2009), 895--901,
\href{https://arxiv.org/abs/0810.3242}{arXiv:0810.3242}.
\bibitem{Ostrik3} V. Ostrik, \emph{Pivotal fusion categories of rank 3},
Moscow Math. J. 15 (2015), 373--396,
\href{https://arxiv.org/abs/1309.4822}{arXiv:1309.4822}.
\bibitem{Deshpande} T. Deshpande, \emph{On centers of bimodule categories
and induction-restriction functors},
\href{https://arxiv.org/abs/1611.04486}{arXiv:1611.04486}.
\bibitem{LPR} Z. Liu, S. Palcoux, Y. Ren, \emph{Classification of Grothendieck
rings of complex fusion categories of multiplicity one up to rank six},
\href{https://arxiv.org/abs/2010.10264}{arXiv:2010.10264}; see the scope
audit for the exact obstruction hypotheses.
\bibitem{Schopieray} A. Schopieray, \emph{Categorification of integral
group rings extended by one dimension}, and the two-Frobenius--Perron-value
results cited in the supplementary notes,
\href{https://arxiv.org/abs/2208.07319}{arXiv:2208.07319}.
\bibitem{Liptrap} J. Liptrap, \emph{Generalized Tambara--Yamagami categories},
\href{https://arxiv.org/abs/1002.3166}{arXiv:1002.3166}.
\bibitem{ENOSolvable} P. Etingof, D. Nikshych, V. Ostrik,
\emph{Weakly group-theoretical and solvable fusion categories},
Adv. Math. 226 (2011), 176--205,
\href{https://arxiv.org/abs/0809.3031}{arXiv:0809.3031}.
\bibitem{NR} D. Nikshych, B. Riepel, \emph{Categorical Lagrangian
Grassmannians and Brauer--Picard groups of pointed fusion categories},
\href{https://arxiv.org/abs/1309.5026}{arXiv:1309.5026}.
\bibitem{HYY} H.-L. Huang, Z. Wan, Y. Ye, \emph{Explicit cocycle formulas
on finite abelian groups with applications to braided linear Gr-categories
and Dijkgraaf--Witten invariants},
\href{https://arxiv.org/abs/1703.03266}{arXiv:1703.03266}.
\bibitem{BN} S. Burciu, S. Natale, \emph{Fusion rules of equivariantizations
of fusion categories},
\href{https://arxiv.org/abs/1206.6625}{arXiv:1206.6625}.
\bibitem{GHR} C. Galindo, S.-M. Hong, E. C. Rowell,
\emph{Generalized and quasi-localizations of braid group representations},
\href{https://arxiv.org/abs/1105.5048}{arXiv:1105.5048},
Theorem 5.20 (unitarity of weakly group-theoretical fusion categories).
\bibitem{NgSchauenburg} S.-H. Ng, P. Schauenburg,
\emph{Frobenius--Schur indicators and exponents of spherical categories},
Adv. Math. 211 (2007), 34--71,
\href{https://arxiv.org/abs/math/0601012}{arXiv:math/0601012}.
\bibitem{EGNO} P. Etingof, S. Gelaki, D. Nikshych, V. Ostrik,
\emph{Tensor Categories}, Mathematical Surveys and Monographs 205,
American Mathematical Society, 2015, Chapters 7--8.
\bibitem{Zesting} C. Delaney, C. Galindo, J. Plavnik, E. C. Rowell, Q. Zhang,
\emph{Braided zesting and its applications}, Commun. Math. Phys. 386
(2021), 1--55, \href{https://arxiv.org/abs/2005.05544}{arXiv:2005.05544}.
\bibitem{Reutter} D. Reutter, \emph{Uniqueness of unitary structure for
unitarizable fusion categories},
\href{https://arxiv.org/abs/1906.09710v2}{arXiv:1906.09710v2},
Theorems~1--2 and Corollary~3.
\bibitem{Izumi} M. Izumi, \emph{A Cuntz algebra approach to the
classification of near-group categories},
\href{https://arxiv.org/abs/1512.04288v2}{arXiv:1512.04288v2},
Theorems~1.1, 4.9 and~6.1, Remark~3.6, Definition~4.8, and Lemma~6.13.
\bibitem{JordanLarson} D. Jordan, E. Larson, \emph{On the classification
of certain fusion categories}, J. Noncommut. Geom. 3 (2009), 481--499,
\href{https://arxiv.org/abs/0812.1603v2}{arXiv:0812.1603v2},
Theorem~1.4, Lemma~2.6, and Proposition~3.1.
\bibitem{ENOHomotopy} P. Etingof, D. Nikshych, V. Ostrik,
with an appendix by E. Meir, \emph{Fusion categories and homotopy theory},
\href{https://arxiv.org/abs/0909.3140v2}{arXiv:0909.3140v2},
Corollary~7.9(i).
\end{thebibliography}
\end{document}
